\documentclass[11pt,reqno]{amsart}

% --- Packages -----------------------------------------------------------------
\usepackage[T1]{fontenc}
\usepackage[utf8]{inputenc}
\usepackage[english]{babel}
\usepackage{amsmath,amssymb,amsthm,mathtools}
\usepackage{microtype}
\usepackage{booktabs}
\usepackage{enumitem}
\usepackage[dvipsnames]{xcolor}
\usepackage{placeins}
\usepackage{hyperref}
\usepackage{xurl}
\hypersetup{
  colorlinks=true,
  linkcolor=blue!50!black,
  citecolor=blue!50!black,
  urlcolor=blue!50!black,
}

\raggedbottom

% --- Theorem environments -----------------------------------------------------
\theoremstyle{plain}
\newtheorem{theorem}{Theorem}[section]
\newtheorem{lemma}[theorem]{Lemma}
\newtheorem{corollary}[theorem]{Corollary}
\newtheorem{proposition}[theorem]{Proposition}
\newtheorem{conjecture}[theorem]{Conjecture}

\theoremstyle{definition}
\newtheorem{definition}[theorem]{Definition}

\theoremstyle{remark}
\newtheorem{remark}[theorem]{Remark}

% --- Macros -------------------------------------------------------------------
\newcommand{\Tab}{T_{a,b}}
\newcommand{\Tneg}{T_{-}}
\newcommand{\Tpos}{T_{+}}
\newcommand{\Z}{\mathbb{Z}}
\newcommand{\N}{\mathbb{N}}
\newcommand{\Q}{\mathbb{Q}}
\newcommand{\Qtwo}{\mathbb{Q}_{2}}
\newcommand{\Ztwo}{\mathbb{Z}_{2}}
\newcommand{\Ztwox}{\mathbb{Z}_{2}^{\times}}
\newcommand{\Nodd}{\N_{\mathrm{odd}}}
\newcommand{\Obs}{\mathcal{O}}
\newcommand{\Atom}{\mathcal{A}}
\newcommand{\Vk}{V_{K}}
\newcommand{\Vi}{V_{I}}
\newcommand{\vk}{v_{K}}
\newcommand{\vi}{v_{I}}
\newcommand{\vtwo}{v_{2}}
\newcommand{\ord}{\mathrm{ord}}
\newcommand{\Phij}{\Phi_{J}}
\newcommand{\Psib}{\Psi_{b}}
\newcommand{\cWa}{c_{W}^{(a)}}
\newcommand{\cWab}{c_{W}^{(a,b)}}
\newcommand{\Obsab}{\Obs_{L}^{(a,b)}}
\newcommand{\Atomab}{\Atom_{L}^{(a,b)}}
\newcommand{\Obstad}{\Obs^{\mathrm{2{\text{-}}ad},(a,b)}}
\newcommand{\srcpath}[1]{\nolinkurl{#1}}
\emergencystretch=3em

% --- Document -----------------------------------------------------------------
\title[Obstruction residues for the $(an+b)$ family]{Obstruction residues for the \texorpdfstring{$(an+b)$}{(an+b)} family:\\
  bias conjugation and multiplier order}

\author{Jakob Stemberger}
\address{Independent researcher, Vienna, Austria}
\email{jakob.stemberger@yaccob.net}
\urladdr{https://orcid.org/0009-0001-2746-6877}
\date{\today}

\subjclass[2020]{11B83, 11S05, 37P05, 37B10, 37B40}

\keywords{Collatz conjecture, $3n\pm1$ map, $(an+b)$ family, residue dynamics,
  parallel reduction, 2-adic density, Wirsching operator}

\begin{document}

\begin{abstract}
For odd integers $a, b$ with $\lvert a \rvert \ge 3$ and $b \ne 0$, the map
$\Tab(n) := (an + b)/2^{\vtwo(an + b)}$ acts on the odd integers. The
\emph{obstruction residues} $\Obsab \subseteq \{1, \ldots, 2^L\}$ are the
odd residues at level $L$ for which an explicit two-track parallel reduction
terminates with a vanishing affine certificate $\Phi_J = 0$; the
$\Tneg$-manuscript proves a constructive lift, positive density $c_W$, and
full factor complexity for the special case $(a,b) = (3,-1)$.

We extend the theory to the whole class via two bridges. Along the bias
axis, the multiplicative conjugation $\Psib(r) := br \bmod 2^L$ is a
shift-preserving bijection $\Obs_L^{(a,1)} \to \Obsab$, so the cardinality
$\lvert \Obsab \rvert$ and the asymptotic density $\cWa$ depend only on
$a$. Along the multiplier axis, a diophantine identity for the reduction
gives, modulo $\lvert a \rvert$, the order lemma
$\delta_J \equiv 0 \pmod{e}$ with $e := \ord_{\lvert a \rvert}(2)$ (the
$\Tneg$-manuscript's parity lemma is the case $a = 3$). The density $\cWa$
exists as a monotone limit, and existence of obstructions ($\cWa > 0$) is
equivalent to a finite diophantine solvability condition; a $2$-adic Haar
interpretation places $\cWa$ in the Bernstein--Lagarias tradition. We
further prove
$L_{\min}(a) \ge e + 1$ (outside the negative-Mersenne case), a
diophantine criterion for the levels carrying
$J = 1$ atoms (with the Mersenne multipliers as its principal solutions),
the no-shift-zero-atom theorem for all $\lvert a \rvert \ge
3$ outside the negative-Mersenne case, and a polynomial bound on fixed-$J$
atom counts with a constructive extremal family at $J = L - 4$ for $a = 3$.
Beyond these theorems the quantitative size of $\cWa$ is so far only
empirical: across the tested multipliers $\cWa > 0$ and
$\cWa \approx C(a)\,4^{-e}$ for an $a$-dependent factor $C(a)$ that stays
bounded in the tested range; whether $C(a)$ is bounded in general
(equivalently $\cWa = \Theta(4^{-e})$) and whether $\alpha(a) < 2$ are the
principal conjectures and remain open (a Wirsching-style operator is
sketched as one route). The
boundary $\lvert a \rvert = 1$ is excluded --- $\Obs_L^{(1,b)} =
\emptyset$ for all $L$ at $a = +1$, while $a = -1$ is degenerate --- confirming
the necessity of $\lvert a \rvert \ge 3$. The negative-Mersenne multipliers
$a = -(2^{v_M} - 1)$ admit a finer sub-class hierarchy treated in a
separate manuscript.
\end{abstract}

\maketitle

% =============================================================================
\section{Introduction}\label{sec:intro}
% =============================================================================

\subsection{Setting}

For odd integers $a, b$ with $\lvert a \rvert \ge 3$ and $b \ne 0$, the
\emph{affine accelerated map} $\Tab$ acts on the set $\Nodd$ of positive
odd integers by
\[
\Tab(n) := \frac{an + b}{2^{\vtwo(an + b)}}.
\]
Two specialisations are classical: $(a,b) = (3,1)$ is the Collatz map
$\Tpos$ (\cite{Lagarias2010,Lagarias2010bib}), and $(a,b) = (3,-1)$ is
the $3n-1$ map $\Tneg$, which is known to admit at least three distinct
cycles (the fixed cycle through $1$, the $\{5,7\}$-cycle, and the
$7$-cycle through $17$; see \cite{Lagarias2010bib}). The Collatz
conjecture asserts that every $\Tpos$-trajectory reaches $1$; the
strongest unconditional progress towards it is Tao's theorem that
almost all orbits attain almost bounded values
\cite{Tao2022}. The analogous question for $\Tneg$ asks whether further
cycles exist beyond the three known ones. Generalised $(an+b)$ maps go back to
Hasse's generalised Syracuse algorithm, whose theory is developed by
Matthews and Watts \cite{MatthewsWatts1984} and M\"oller
\cite{Moeller1978} (the latter's title is literally ``Hasse's
generalisation of the Syracuse algorithm''). The broader class of
piecewise-defined generalised Collatz iterations is computationally
universal --- Conway's undecidability of such iteration problems
\cite{Conway1972} --- which is part of what makes an exact,
level-by-level algebraic characterisation of a structured subfamily
worthwhile; the present obstruction-residue viewpoint is exactly such a
characterisation, rather than a measure-theoretic or heuristic one. The
2-adic conjugation that underlies our bias-axis bridge
(Section~\ref{sec:tools}) and the Haar interpretation of
Section~\ref{sec:adic} are
in the tradition of the $3x+1$ conjugacy map of Bernstein and Lagarias
\cite{BernsteinLagarias1996}, who realise the Collatz dynamics as a
conjugate of the shift on $\Ztwo$.

Both maps share an algebraic feature that turns out to be visible at the
level of \emph{residues modulo $2^L$}. The $\Tneg$-manuscript
\cite{ObstructionResidues2026} introduces
a family of residue classes $\Obs_L \subseteq \{1, \ldots, 2^L\}$
characterised by an algebraic criterion on a two-track parallel
reduction, and proves three structural results for this family:
\begin{enumerate}[label=(\arabic*)]
\item a constructive lift between modular levels $L$ and $L+1$;
\item a strictly positive asymptotic density $c_W \in (0,\tfrac12]$ (only
  odd residues can be obstructions, so $\tfrac12$ is the trivial ceiling)
  with the rigorous lower bound $c_W \ge 7556/65536 \approx 0.115$;
\item full factor complexity $p_W(\ell) = 2^\ell$ of the associated
  language of binary representations.
\end{enumerate}
All three rest on a \emph{parity lemma}: at termination of the parallel
reduction, the difference $\delta_J = \Vk^{(J)} - \Vi^{(J)} - v$ of
cumulative valuations is forced to be an even integer.

\subsection{The two bridges}

The present paper develops the same theory in the wider $(an+b)$ setting,
$\lvert a \rvert \ge 3$ and $b$ odd, by isolating two bridges from the
$\Tneg$-manuscript's special case.

\medskip
\noindent\textbf{Bridge 1: bias axis (multiplicative conjugation).} The
multiplicative map $\Psib(r) := br \bmod 2^L$ is a shift-preserving
bijection $\Obs_L^{(a,1)} \to \Obsab$ for every odd $b \ne 0$
(Theorem~\ref{thm:conjugation}). The $\Tneg$-manuscript's involution $r \mapsto
2^L - r$ that transfers results from $\Tneg$ to $\Tpos$
(\cite[Section~7]{ObstructionResidues2026}) is the special case $\Psi_{-1}$.
As a consequence, the cardinalities $\lvert\Obsab\rvert$ and the
asymptotic density $\cWab$ depend only on $a$, not on $b$
(Corollary~\ref{cor:b-universality}). We write $\cWa := \cWab$
unambiguously.

\medskip
\noindent\textbf{Bridge 2: multiplier axis (order lemma).} Reducing the
\emph{diophantine identity} for $\Phi_J = 0$ modulo $\lvert a \rvert$
forces $\delta_J \equiv 0 \pmod{e}$, where $e := \ord_{\lvert a \rvert}(2)$
(Lemma~\ref{lem:order-lemma}). The $\Tneg$-manuscript's parity lemma is the
case $a = 3, e = 2$. Unlike the conjugation theorem, the order lemma
restricts but does not preserve structure across $a$: the step size $e$
of the shift index grows with the multiplicative order and drives an
empirical density law $\cWa \approx C(a)\,4^{-e}$ across the tested range
$\lvert a \rvert \in \{3, 5, 7, 9, 15, 17, 21, 31\}$
(Section~\ref{sec:scaling}).

\subsection{Main results}

For odd $a, b$ with $\lvert a \rvert \ge 3$ and $b \ne 0$ and a level $L
\ge 1$, the obstruction set $\Obsab$ at level $L$ is defined in
Section~\ref{sec:setup} (Definition~\ref{def:obstruction}) as the set
of odd residues $r$ for which the two-track parallel reduction
terminates with $\Phi_J = 0$. The shift index $\delta_J = \delta_J(r,L)$
is the difference of cumulative valuations at termination; the constant
$e := \ord_{\lvert a \rvert}(2)$ is the multiplicative order of $2$
modulo $\lvert a \rvert$.

\begin{theorem}[Conjugation, b-axis universality]\label{thm:conjugation-intro}
For every odd $b \ne 0$ the multiplicative map $\Psib(r) := br \bmod 2^L$
is a shift-preserving bijection $\Obs_L^{(a,1)} \to \Obsab$. In
particular, the cardinality $\lvert\Obsab\rvert$ and the asymptotic
density $\cWab = \lim_L \lvert\Obsab\rvert/2^L$ depend only on $a$.
\end{theorem}

\begin{theorem}[Diophantine identity, order lemma]\label{thm:order-intro}
For every $r \in \Obsab$ with termination index $J$,
\[
a^{J+1} = \sum_{j=0}^{J} a^{J-j} X_j, \qquad X_j := 2^{\Vk^{(j)}} - 2^{\Vi^{(j)} + v}.
\]
Reducing modulo $\lvert a \rvert$ gives $\delta_J \equiv 0 \pmod{e}$.
\end{theorem}

\begin{theorem}[Density limit and existence criterion]\label{thm:adic-intro}
The density $\cWab = \lim_L \lvert \Obsab \rvert / 2^L$ exists as a
non-decreasing limit, and $\cWa > 0$ if and only if the diophantine
identity above is realisable for at least one finite valuation tuple
$(J, \Vk^{(j)}, \Vi^{(j)}, v)$.
\end{theorem}

\noindent A measure-theoretic reading of this density --- $\cWa =
\mu_2(\Obstad)$, the Haar measure of the open set $\Obstad \subseteq
\Ztwox$ of $2$-adic units that are eventually obstructions --- is recorded
as Remark~\ref{rem:haar}, in the Bernstein--Lagarias tradition.

\begin{theorem}[No shift-zero atom, generalised]\label{thm:no-G0-intro}
Assume $a$ is not a negative Mersenne number $-(2^{v_M} - 1)$. For every
level $L$,
\[
\lvert \Atom_L^{(a,b), G_0} \rvert = 0;
\]
every atomic obstruction has nonzero shift index.
\end{theorem}

\begin{proposition}[Atom localisation at fixed $J$]\label{prop:atom-J-intro}
At fixed termination index $J$ the atomic $G_{\ne 0}$-count is bounded by
$\lvert \Atom_L^{(a,b), G_{\ne 0}} \rvert_J \le \binom{L-1}{J} = O_J(L^{J})$
(Proposition~\ref{prop:atom-J}): every such atom drives one track to the
boundary of its modulus, and the boundary valuation data then determine the
residue uniquely. For $(a,b) = (3,-1)$ the residue
$r := 2^{L-2} + 3$ is an atomic obstruction at every \emph{even} level
$L \ge 6$ with $J = L - 4$ and shift index $1$, realising the empirical
extreme $J_{\max} = L - 4$.
\end{proposition}

\begin{theorem}[Degenerate edge case]\label{thm:edge-intro}
For $a = +1$ and every odd $b$, $\Obsab = \emptyset$ for every $L$. For
$a = -1$ the map degenerates ($an + b$ vanishes on part of the domain)
and direct enumeration gives $\lvert \Obsab \rvert = 2^{L-2}$
(Remark~\ref{rem:a-minus-1}); neither value of $\lvert a \rvert
= 1$ supports the genuine obstruction dynamics that require $\lvert a
\rvert \ge 3$.
\end{theorem}

Theorems~\ref{thm:conjugation-intro}--\ref{thm:no-G0-intro},
Proposition~\ref{prop:atom-J-intro} and Theorem~\ref{thm:edge-intro} are proven
across Sections~\ref{sec:tools} (conjugation and order lemma),
\ref{sec:adic} (density limit and existence criterion),
\ref{sec:structure} (no-shift-zero atom
and atom localisation), and \ref{sec:edge} (degenerate edge case). Two further results
(Theorem~\ref{thm:Lmin-lower} on the lower bound for the smallest
non-trivial level and Proposition~\ref{prop:J1-character} on the
characterisation of $J = 1$ atoms, for positive $a$) are stated in
Section~\ref{sec:Lmin}.
The principal open problems --- $\alpha(a) < 2$ and the existence
conjecture $\cWa > 0$ for non-Mersenne
$\lvert a \rvert$ --- are discussed in Section~\ref{sec:spectral}, and the
empirical density scaling $\cWa \approx C(a)\,4^{-e}$ (with the
boundedness of $C(a)$ itself conjectural) is reported in
Section~\ref{sec:scaling}.

\subsection{Relation to the negative-Mersenne case}

For $a = -(2^{v_M} - 1)$ Mersenne with $v_M \ge 2$, the $(an+b)$ theory
admits a finer sub-class hierarchy indexed by the shift index $j \in
\Z$ (the iterated $A$-, $B$-, $C$-, and $F_1$-families, plus a
$j = -1$ family $G_1$ specific to negative $a$); the $\Tneg$-manuscript-style
density bound becomes $\cWa \ge 1/(2\lvert a \rvert)$, asymptotically
sharp in $v_M$. We treat this case in a separate manuscript (in
preparation) and only cite the relationship where it sharpens the general
theory.

\subsection{Outline}

Section~\ref{sec:setup} fixes notation and defines the two-track parallel
reduction. Section~\ref{sec:tools} proves the four foundational tools
(conjugation, order lemma, diophantine identity, and the series
representation of $\cWa$). Section~\ref{sec:adic} establishes the density
limit and reduces existence of obstructions to a finite diophantine
solvability condition (with a $2$-adic Haar interpretation).
Section~\ref{sec:Lmin} establishes the lower bound $L_{\min}(a) \ge e + 1$
and characterises the levels at which $J = 1$ atoms exist.
Section~\ref{sec:structure} contains the structural results over $a$:
the generalised no-shift-zero-atom theorem, polynomial atom bounds at
fixed $J$, and the constructive $L - 4$ extremal family.
Section~\ref{sec:spectral} discusses the principal open problem
$\alpha(a) < 2$ (with a Wirsching-style operator sketched as one route).
Section~\ref{sec:scaling} reports the empirical scaling $\cWa = \Theta(4^{-e})$.
Section~\ref{sec:edge} treats the degenerate boundary $\lvert a \rvert = 1$
($a = +1$ empty, $a = -1$ degenerate).
Section~\ref{sec:neg-mersenne} records the relationship to the companion
manuscript's negative-Mersenne sub-class analysis.
Appendix~\ref{sec:errata} lists two
refuted hypotheses that motivated the present development.

% =============================================================================
\section{Setup: parallel reduction and obstruction residues}\label{sec:setup}
% =============================================================================

\subsection{The two-track parallel reduction}

Fix odd integers $a, b$ with $\lvert a \rvert \ge 3$ and $b \ne 0$. Let
$r$ be a positive odd integer, and write
\[
v := \vtwo(r + b), \qquad m := (r + b)/2^v.
\]
Since $b$ is odd (which we always assume) and $r$ is odd, the sum
$r + b$ is even, so $v = \vtwo(r+b) \ge 1$ and $m = (r+b)/2^v$ is odd.

Fix a level $L > v$. The parallel reduction proceeds simultaneously on
two tracks:
\begin{itemize}
\item the \emph{K-track}, $(a_K^{(0)}, b_K^{(0)}) := (r, 2^L)$;
\item the \emph{I-track}, $(a_I^{(0)}, b_I^{(0)}) := (m, 2^{L-v})$.
\end{itemize}
At each step $j \ge 0$, let
$A_K^{(j)} := a \cdot a_K^{(j)} + b$, $A_I^{(j)} := a \cdot a_I^{(j)} + b$,
$\vk^{(j)} := \vtwo(A_K^{(j)})$, $\vi^{(j)} := \vtwo(A_I^{(j)})$, and
update
\[
a_K^{(j+1)} := A_K^{(j)}/2^{\vk^{(j)}}, \qquad b_K^{(j+1)} := a \cdot b_K^{(j)}/2^{\vk^{(j)}},
\]
provided $\vk^{(j)} < \vtwo(b_K^{(j)})$; otherwise the K-track terminates
at index $j$. The I-track is updated analogously with $\vi^{(j)}$ and
$b_I^{(j)}$. Cumulative valuations are
$\Vk^{(j)} := \sum_{i<j} \vk^{(i)}$ and
$\Vi^{(j)} := \sum_{i<j} \vi^{(i)}$, so that the residual valuations of
the moduli are $L - \Vk^{(j)}$ and $L - v - \Vi^{(j)}$ respectively.

The reduction proceeds as long as both tracks can step; it terminates as
soon as either condition fails. We write $J = J(r,L)$ for the index of
termination.

\subsection{The affine certificate}

\begin{lemma}\label{lem:affine}
For each $j$ with $0 \le j \le J$, there exist $c_j, d_j \in \Q$ with
\[
a_I^{(j)} = c_j \cdot a_K^{(j)} + d_j \qquad \text{in } \Q,
\]
where $c_0 = 2^{-v}$, $d_0 = b \cdot 2^{-v}$, and
\[
c_{j+1} = c_j \cdot 2^{\vk^{(j)} - \vi^{(j)}},
\qquad
d_{j+1} = \frac{a \cdot d_j + b(1 - c_j)}{2^{\vi^{(j)}}}.
\]
\end{lemma}

\begin{proof}
At $j = 0$: $m = (r+b)/2^v = 2^{-v} r + b \cdot 2^{-v}$. Inductively,
$A_I^{(j)} = a(c_j a_K^{(j)} + d_j) + b = c_j A_K^{(j)} + (a d_j + b(1 - c_j))$.
Dividing by $2^{\vi^{(j)}}$ and using $A_K^{(j)} = 2^{\vk^{(j)}} a_K^{(j+1)}$:
\[
a_I^{(j+1)} = c_j \cdot 2^{\vk^{(j)} - \vi^{(j)}} \cdot a_K^{(j+1)}
            + \frac{a d_j + b(1 - c_j)}{2^{\vi^{(j)}}}. \qedhere
\]
\end{proof}

Write $\delta_j := \Vk^{(j)} - \Vi^{(j)} - v$. From the update rule,
$c_j = 2^{\delta_j}$ for every $j$. Define the \emph{affine certificate}
at termination as
\[
\Phi_J := b(1 - c_J) + a \cdot d_J.
\]

\begin{definition}[Obstruction residue]\label{def:obstruction}
Let $r$ be a positive odd integer with $r + b$ even, $v := \vtwo(r+b) \ge 1$,
and $L > v$. We call $r$ an \emph{obstruction residue for $\Tab$ at
level $L$} if $\Phi_J(r, L) = 0$. The set of all such residues in
$\{1, \ldots, 2^L\}$ is denoted $\Obsab$.
\end{definition}

\begin{remark}[Terminal synchronisation]\label{rem:sync}
The defining condition $\Phi_J = 0$ is equivalent to an \emph{exact}
synchronisation of the two tracks at termination. From the affine relation
$a_I^{(j)} = c_j a_K^{(j)} + d_j$ and $A_K^{(j)} = a\,a_K^{(j)} + b$ one has
the identity
\[
A_I^{(j)} = a\,a_I^{(j)} + b = c_j A_K^{(j)} + \Phi_j, \qquad
\Phi_j := b(1 - c_j) + a\,d_j,
\]
so $\Phi_J = 0 \iff A_I^{(J)} = c_J A_K^{(J)}$. Since
$\vtwo(A_K^{(J)}) = \vk^{(J)}$ and $c_J = 2^{\delta_J}$, this forces
$\vi^{(J)} = \vk^{(J)} + \delta_J$; substituting this and
$\Vi^{(J)} = \Vk^{(J)} - v - \delta_J$ into the two stop conditions
$\vk^{(J)} \ge L - \Vk^{(J)}$ and $\vi^{(J)} \ge L - v - \Vi^{(J)}$ makes
them identical, so \emph{both tracks meet the termination condition at the
same index $J$}. We call the terminal step \emph{exact} when
$\vk^{(J)} = L - \Vk^{(J)}$ (equivalently the I-track equality holds) and
\emph{strict} when $\vk^{(J)} > L - \Vk^{(J)}$; by the synchronisation the
two tracks share this type. The certificate is thus the algebraic shadow
of a synchronisation of stopping times for $\Tab$-iterates of $r$ and $m$;
the present paper develops only the algebraic side.
\end{remark}

\subsection{The shift index}

The integer $\delta_J(r, L) = \Vk^{(J)} - \Vi^{(J)} - v$ controls the
fine structure of $\Obsab$. We refer to it as the \emph{shift index}.
The shift-zero set is
\[
\Obs_L^{(a,b), G_0} := \{r \in \Obsab : \delta_J(r,L) = 0\},
\qquad
\Obs_L^{(a,b), G_{\ne 0}} := \Obsab \setminus \Obs_L^{(a,b), G_0},
\]
and analogously for $\Atom_L^{(a,b), G_0}, \Atom_L^{(a,b), G_{\ne 0}}$
once atomic obstructions are defined below.

\subsection{Lifts and atoms}\label{ssec:lifts}

The obstruction sets at consecutive levels are linked by a lift map,
which we record here because the density series
(Corollary~\ref{cor:series}), the density limit
(Section~\ref{sec:adic}), and the lower bound of
Section~\ref{sec:scaling} all rest on it.

\begin{definition}[Lift, atom]\label{def:atom}
Let $r \in \Obsab$. The two residues $r$ and $r + 2^L$ in
$\{1, \ldots, 2^{L+1}\}$ are the \emph{level-$(L+1)$ lifts} of $r$. An
obstruction $r \in \Obsab$ is \emph{atomic} if it is not a lift of any
level-$(L-1)$ obstruction, i.e.\ $r \bmod 2^{L-1} \notin
\Obs_{L-1}^{(a,b)}$; the set of atomic obstructions at level $L$ is
denoted $\Atomab$, and $\Atom_L^{(a,b), G_0}$, $\Atom_L^{(a,b), G_{\ne 0}}$
are its shift-zero and nonzero-shift parts.
\end{definition}

\begin{theorem}[Lift theorem]\label{thm:lift}
For every $r \in \Obsab$ both level-$(L+1)$ lifts $r$ and $r + 2^L$ lie in
$\Obs_{L+1}^{(a,b)}$, and conversely every non-atomic obstruction at level
$L+1$ is a lift of a unique obstruction at level $L$. Consequently, with
$h_L := \lvert \Obsab \rvert$,
\[
h_{L+1} = 2 h_L + \lvert \Atom_{L+1}^{(a,b)} \rvert.
\]
\end{theorem}

\begin{proof}
For $(a,b) = (3,-1)$ this is the constructive lift of
\cite[Theorems~3.1 and~3.2]{ObstructionResidues2026}. We recall the mechanism and
its $a$-dependence. The lift $r_+ := r$ runs the same reduction with the
modulus doubled, so its data agree with those of $r$ at level $L$ until
the level-$(L+1)$ stop. By the terminal synchronisation
(Remark~\ref{rem:sync}) the two tracks share an exact/strict terminal type,
and the two cases are: if the terminal step at $L$ was \emph{exact}, one
extra step at the terminal index yields, using
$\vk^{(J)} - \vi^{(J)} = -\delta_J$,
$c_{J+1} = c_J \cdot 2^{\vk^{(J)} - \vi^{(J)}} = 1$ and $d_{J+1} = 0$ (a
shift-zero lift); if it was \emph{strict}, the index-$J$ data already satisfy
the level-$(L+1)$ stop and the terminal certificate is unchanged (a
shift-$s$ lift); both are obstructions. The
lift $r_- := r + 2^L$ differs from $r_+$ by the discrepancy
$\Delta_K^{(j)} = 2^{L - \Vk^{(j)}} a^j$ (and $\Delta_I^{(j)} =
2^{L - v - \Vi^{(j)}} a^j$); because $a$ is odd, multiplying by $a$
preserves the $2$-adic valuation, so while $\Vk^{(j)} < L$ the two
summands of $a\,a_{K,-}^{(j)} + b = (a\,a_{K,+}^{(j)} + b) + a\,\Delta_K^{(j)}$
sit at distinct valuations and the ultrametric identity forces the
per-step valuations to match; the case split on the exact/strict terminal
type is identical and yields the second lift. The multiplier $a$ enters only
through $c_J = 2^{\delta_J}$, $d_J = b(c_J - 1)/a$ (replacing the $4^{s}$
and $(1 - 4^{s})/3$ of the $a = 3$ case) and the factor $a^j$ in the
discrepancy --- none of which affects the lift count of two. The argument
therefore holds for every odd $a$ with $\lvert a \rvert \ge 3$ and odd
$b \ne 0$. For the converse and the balance: an obstruction
$r \in \Obs_{L+1}^{(a,b)}$ is either atomic --- contributing to
$\lvert \Atom_{L+1}^{(a,b)} \rvert$ --- or, by Definition~\ref{def:atom},
satisfies $r \bmod 2^{L} \in \Obsab$, in which case it is one of the two
lifts of its unique parent $r \bmod 2^{L}$; the two alternatives are mutually
exclusive and exhaustive, so $h_{L+1} = 2 h_L + \lvert \Atom_{L+1}^{(a,b)}
\rvert$. (Direct enumeration confirms the balance for all tested $(a,b)$ and
$L \le 16$.) \qedhere
\end{proof}

\subsection{Asymptotic density}

The \emph{asymptotic obstruction density} is
\[
\cWab := \lim_{L \to \infty} \frac{\lvert \Obsab \rvert}{2^L}
\]
when this limit exists. The $\Tneg$-manuscript proves existence and positivity
of $c_W = c_W^{(3,-1)}$; we extend both to the entire $(an+b)$ family in
Section~\ref{sec:tools} (cardinality preservation under $\Psib$) and
Section~\ref{sec:adic} (existence of the limit via the lift theorem).

% =============================================================================
\section{Foundational tools}\label{sec:tools}
% =============================================================================

\subsection{The conjugation theorem}

\begin{theorem}[Conjugation theorem]\label{thm:conjugation}
Let $a$ be an odd integer with $\lvert a \rvert \ge 3$ and $b$ a nonzero
odd integer. For every level $L \ge 1$ the multiplicative map
\[
\Psib : \Z/2^L\Z \to \Z/2^L\Z, \qquad \Psib(r) := b r \bmod 2^L,
\]
restricts to a bijection $\Obs_L^{(a,1)} \to \Obsab$ that preserves the
shift index: $\delta_J^{(a,b)}(\Psib(r)) = \delta_J^{(a,1)}(r)$ (trivially
so at levels below $L_{\min}(a)$, where both sets are empty).
\end{theorem}

\begin{proof}
Since $b$ is odd, $\Psib$ is a bijection of $\Z/2^L\Z$. We must show it
restricts to obstructions and preserves the shift index.

\emph{Step 1 (initial data).} For $r \in \Obs_L^{(a,1)}$,
$\vtwo(r+1) = v$ and $m_K := (r+1)/2^v$ is odd. For $r' := \Psib(r) = br$,
\[
r' + b = br + b = b(r+1), \qquad \vtwo(r'+b) = \vtwo(b) + \vtwo(r+1) = 0 + v = v,
\]
since $b$ is odd. So $v_K^{\prime,(0)} := \vtwo(r' + b)$ equals $v$, and
$m_K' := (r'+b)/2^v = b \cdot m_K$. The initial affine data satisfy
$c_0^{(a,b)} = 2^{-v} = c_0^{(a,1)}$ and
$d_0^{(a,b)} = b \cdot 2^{-v} = b \cdot d_0^{(a,1)}$.

\emph{Step 2 (step-by-step conjugation).} We claim by induction:
\begin{equation}\label{eq:conj-induction}
a_K^{(a,b),(j)} \equiv b \cdot a_K^{(a,1),(j)} \pmod{2^{L - \Vk^{(j)}}},
\qquad
v_K^{(a,b),(j)} = v_K^{(a,1),(j)},
\end{equation}
and analogously for the I-track. Base case $j = 0$: $a_K^{(a,b),(0)} = b r
= b \cdot a_K^{(a,1),(0)}$, and we showed $v_K$ matches at step $0$.
Inductive step: write $a_K^{(a,b),(j)} = b \cdot a_K^{(a,1),(j)} + 2^{L -
\Vk^{(j)}} \cdot \eta_j$ for some integer $\eta_j$. Then
\[
A_K^{(a,b),(j)} = a (b a_K^{(a,1),(j)} + 2^{L - \Vk^{(j)}} \eta_j) + b
                = b A_K^{(a,1),(j)} + 2^{L - \Vk^{(j)}} \cdot a \eta_j.
\]
Since $b$ is odd, $\vtwo(b A_K^{(a,1),(j)}) = \vtwo(A_K^{(a,1),(j)}) =
v_K^{(a,1),(j)}$; since $v_K^{(a,1),(j)} < L - \Vk^{(j)}$
(non-termination), the second term has higher 2-adic valuation, so
$v_K^{(a,b),(j)} = v_K^{(a,1),(j)}$ and dividing through gives
$a_K^{(a,b),(j+1)} \equiv b a_K^{(a,1),(j+1)} \pmod{2^{L - \Vk^{(j+1)}}}$.

\emph{Step 3 (termination and certificate).} The matching valuations
yield $\Vk^{(a,b),(j)} = \Vk^{(a,1),(j)}$ and
$\Vi^{(a,b),(j)} = \Vi^{(a,1),(j)}$ for every $j$, so termination
occurs at the same index $J$ for both. From the update rule for $(c_j,
d_j)$ and the initial data of Step~1, induction gives
$c_J^{(a,b)} = c_J^{(a,1)}$ and $d_J^{(a,b)} = b \cdot d_J^{(a,1)}$.
Therefore
\[
\Phi_J^{(a,b)} = b(1 - c_J^{(a,b)}) + a \cdot d_J^{(a,b)}
              = b \bigl[(1 - c_J^{(a,1)}) + a \cdot d_J^{(a,1)}\bigr]
              = b \cdot \Phi_J^{(a,1)},
\]
and since $b \ne 0$, $\Phi_J^{(a,b)} = 0 \iff \Phi_J^{(a,1)} = 0$.

\emph{Step 4 (shift index).} $\delta_J = \Vk^{(J)} - \Vi^{(J)} - v$
is unchanged because all summands are unchanged. \qedhere
\end{proof}

\begin{corollary}[$\Tneg$ to $\Tpos$ involution as $\Psi_{-1}$]\label{cor:involution}
The $\Tneg$-manuscript's involution $r \mapsto 2^L - r$ inducing a bijection
$\Obs_L^{(3,-1)} \leftrightarrow \Obs_L^{(3,1)}$
(\cite[Section~7]{ObstructionResidues2026}) is the special case $\Psi_{-1}$ of
Theorem~\ref{thm:conjugation}.
\end{corollary}

\begin{corollary}[b-axis universality of cardinality and density]\label{cor:b-universality}
For every odd $b \ne 0$ and every $L \ge 1$,
\[
\lvert \Obsab \rvert = \lvert \Obs_L^{(a,1)} \rvert.
\]
In particular, $\cWab = \cWa$ depends only on $a$.
\end{corollary}

We have verified Theorem~\ref{thm:conjugation} numerically for
$a = 3$, $L = 12$, and all odd $b \in \{-7, -5, \ldots, +13\}$: the
cardinalities $\lvert \Obs_{12}^{(3, b)} \rvert = 409$ are identical and
the shift-index assignments match under $\Psib$
(code: \srcpath{code/python/factor_lang_check.py}).

\subsection{The diophantine identity}

\begin{theorem}[Diophantine identity]\label{thm:diophantine}
For every $r \in \Obsab$ with termination index $J$, set
\[
X_j := 2^{\Vk^{(j)}} - 2^{\Vi^{(j)} + v}, \qquad 0 \le j \le J.
\]
Then
\begin{equation}\label{eq:diophantine}
a^{J+1} = \sum_{j=0}^{J} a^{J-j} \cdot X_j.
\end{equation}
\end{theorem}

\begin{proof}
Let $D_j := 2^{\Vi^{(j)} + v} \cdot d_j$. From the $d$-update of
Lemma~\ref{lem:affine},
\begin{align*}
D_{j+1} = 2^{\Vi^{(j+1)} + v} d_{j+1}
  &= 2^{\Vi^{(j)} + \vi^{(j)} + v} \cdot \frac{a d_j + b(1 - c_j)}{2^{\vi^{(j)}}} \\
  &= a D_j + b \cdot 2^{\Vi^{(j)} + v}(1 - c_j).
\end{align*}
Since $c_j = 2^{\delta_j} = 2^{\Vk^{(j)} - \Vi^{(j)} - v}$,
$2^{\Vi^{(j)} + v}(1 - c_j) = 2^{\Vi^{(j)} + v} - 2^{\Vk^{(j)}} = -X_j$,
so the recurrence simplifies to
\begin{equation}\label{eq:D-recurrence}
D_{j+1} = a D_j - b \cdot X_j, \qquad D_0 = b.
\end{equation}
Unrolling,
\[
D_{J+1} = a^{J+1} D_0 - b \sum_{j=0}^{J} a^{J-j} X_j
       = b \biggl( a^{J+1} - \sum_{j=0}^{J} a^{J-j} X_j \biggr).
\]
On the other hand, the certificate $\Phi_J = b(1 - c_J) + a d_J = 0$
multiplied by $2^{\Vi^{(J)} + v}$ gives $a D_J = b \cdot X_J$, and one
step of \eqref{eq:D-recurrence} gives $D_{J+1} = a D_J - b X_J = 0$.
Combining the two expressions for $D_{J+1}$ and using $b \ne 0$,
\eqref{eq:diophantine} follows. \qedhere
\end{proof}

\begin{remark}\label{rem:b-independence}
The diophantine identity~\eqref{eq:diophantine} is independent of $b$: it
involves only $a$, the cumulative valuations $\Vk^{(j)}, \Vi^{(j)}$, and
$v$. This is the structural reason behind the conjugation theorem:
$\Psib$ preserves the entire valuation sequence (Step~3 of the proof of
Theorem~\ref{thm:conjugation}), so an obstruction at $(a,1)$ produces
the same valuation tuple at $(a,b)$ and~\eqref{eq:diophantine} is
satisfied for the conjugated residue as well.
\end{remark}

\subsection{The order lemma}

\begin{lemma}[Order lemma]\label{lem:order-lemma}
Let $e := \ord_{\lvert a \rvert}(2)$, the multiplicative order of $2$
modulo $\lvert a \rvert$. For every $r \in \Obsab$ with termination
index $J$,
\[
\delta_J \equiv 0 \pmod{e}.
\]
\end{lemma}

\begin{proof}
Reducing~\eqref{eq:diophantine} modulo $\lvert a \rvert$ leaves only the
$j = J$ term:
\[
0 \equiv X_J = 2^{\Vk^{(J)}} - 2^{\Vi^{(J)} + v} \pmod{\lvert a \rvert}.
\]
Since $\gcd(2, \lvert a \rvert) = 1$, the two powers are congruent and
both invertible modulo $\lvert a \rvert$, so
$2^{\Vk^{(J)} - (\Vi^{(J)} + v)} = 2^{\delta_J} \equiv 1
\pmod{\lvert a \rvert}$. By definition of $e$ this gives $e \mid
\delta_J$ when $\delta_J \ge 0$; when $\delta_J < 0$ the same congruence
reads $2^{-\delta_J} \equiv 1 \pmod{\lvert a \rvert}$, so $e \mid
(-\delta_J)$, hence $e \mid \delta_J$ in either case. \qedhere
\end{proof}

For $a = 3$, $e = \ord_3(2) = 2$, and Lemma~\ref{lem:order-lemma}
specialises to the $\Tneg$-manuscript's parity lemma
(\cite[Lemma~2.4]{ObstructionResidues2026}): $\delta_J$ is even. The
\emph{shift-$s$ normal form} of the $\Tneg$-manuscript becomes the
\emph{shift-$s$ normal form to step $e$}, $\delta_J = s \cdot e$ with
$s \in \Z$. Numerical values of $e$ for small $\lvert a \rvert$ are
collected in Table~\ref{tab:order}.

\begin{table}[ht]
\centering
\begin{tabular}{rrl}
\toprule
$\lvert a \rvert$ & $e = \ord_{\lvert a\rvert}(2)$ & shift step \\
\midrule
3 & 2 & $\delta_J \in 2\Z$ ($\Tneg$ parity) \\
5 & 4 & $\delta_J \in 4\Z$ \\
7 & 3 & $\delta_J \in 3\Z$ \\
9 & 6 & $\delta_J \in 6\Z$ \\
11 & 10 & $\delta_J \in 10\Z$ \\
13 & 12 & $\delta_J \in 12\Z$ \\
15 & 4 & $\delta_J \in 4\Z$ \\
17 & 8 & $\delta_J \in 8\Z$ \\
31 & 5 & $\delta_J \in 5\Z$ \\
\bottomrule
\end{tabular}
\caption{Multiplicative order $e = \ord_{\lvert a\rvert}(2)$ and the
induced shift-index step size, for small odd $\lvert a \rvert$.}
\label{tab:order}
\end{table}

\subsection{Series representation of the density}

The diophantine identity also produces a clean series for $\cWa$. Let
$h_L := \lvert \Obsab \rvert$ and $a_L := \lvert \Atom_L^{(a,b),
G_{\ne 0}} \rvert$ (atomic count of nonzero-shift obstructions at level
$L$).

\begin{corollary}[Series representation]\label{cor:series}
By the lift theorem (Theorem~\ref{thm:lift}), $h_{L+1} = 2 h_L +
a_{L+1} + \varepsilon_L$, where $\varepsilon_L := \lvert \Atom_{L+1}^{(a,b),
G_0} \rvert$ tracks the $G_0$-atomic contributions. Telescoping this
balance gives
\[
\cWa = \sum_{l \ge L_{\min}(a)} \frac{a_l}{2^l} + c_W^{(a), G_0},
\]
where $c_W^{(a), G_0} := \sum_l \varepsilon_{l-1}/2^l$ is the contribution
of the shift-zero \emph{atoms} (not of the shift-zero obstructions, which
are lifts already counted in the $2 h_L$ term and in fact carry the bulk
of the density). By Theorem~\ref{thm:no-G0-atom-general} the summand
$\varepsilon_{l-1}$ vanishes at every level (outside the negative-Mersenne
case), so $c_W^{(a), G_0} = 0$; this recovers the $\Tneg$-manuscript's
$c_W^{(3), G_0} = 0$
(\cite[Theorem~6.2]{ObstructionResidues2026}) as the case $a = 3$.
\end{corollary}

The proof is a telescoping of the lift balance and is identical to the
$\Tneg$-manuscript's (\cite[Section~6.3]{ObstructionResidues2026}); the
$\Tneg$-specific input is only the no-$G_0$-atom statement, generalised
to all $\lvert a \rvert \ge 3$ outside the negative-Mersenne case in
Theorem~\ref{thm:no-G0-atom-general}.

\subsection{Determinacy from valuation data}

A fixed sequence of step valuations pins the seed to a single residue
class. We isolate this as a lemma, used both by the existence criterion
(Section~\ref{sec:adic}) and by the atom bound
(Proposition~\ref{prop:atom-J}).

\begin{lemma}[Determinacy]\label{lem:determinacy}
Fix an integer $v \ge 1$, a termination index $J \ge 0$, and step
valuations $\vk^{(0)}, \ldots, \vk^{(J-1)} \ge 1$ and
$\vi^{(0)}, \ldots, \vi^{(J-1)} \ge 1$; write
$\Vk^{(j)} := \sum_{i<j}\vk^{(i)}$ and $\Vi^{(j)} := \sum_{i<j}\vi^{(i)}$.
The set of odd seeds $r$ whose parallel reduction realises exactly this
valuation sequence is either empty or a single residue class. Fixing the
K-data alone pins $r$ modulo $2^{\Vk^{(J)}+1}$; fixing $v$ together with the
I-data pins $r$ modulo $2^{\Vi^{(J)}+v+1}$.
\end{lemma}

\begin{proof}
By induction the K-iterate is affine in the seed $r$:
\[
a_K^{(j)} = \frac{a^j r + \beta_j}{2^{\Vk^{(j)}}}, \qquad
\beta_0 = 0, \quad \beta_{j+1} = a\beta_j + b\,2^{\Vk^{(j)}},
\]
the recursion for $\beta_j$ following from
$a_K^{(j+1)} = (a\,a_K^{(j)} + b)/2^{\vk^{(j)}}$. The requirement
$\vtwo(a\,a_K^{(j)} + b) = \vk^{(j)}$ reads
$a\,a_K^{(j)} + b \equiv 2^{\vk^{(j)}} \pmod{2^{\vk^{(j)}+1}}$; clearing the
denominator $2^{\Vk^{(j)}}$ and substituting the affine form turns it into
\[
a^{j+1} r \equiv 2^{\Vk^{(j+1)}} - a\beta_j - b\,2^{\Vk^{(j)}}
            \pmod{2^{\Vk^{(j+1)}+1}}.
\]
Since $a$ is odd, $a^{j+1}$ is a unit modulo $2^{\Vk^{(j+1)}+1}$, so this
fixes $r$ modulo $2^{\Vk^{(j+1)}+1}$. The conjunction over
$j = 0, \ldots, J-1$ has nested moduli, so the set of seeds realising the
entire K-sequence is empty or a single residue class modulo
$2^{\Vk^{(J)}+1}$. The identical argument on the I-track --- whose seed
$m = (r+b)/2^v$ is affine in $r$ once $v := \vtwo(r+b)$ is fixed --- shows
that fixing $v$ together with $\vi^{(0)}, \ldots, \vi^{(J-1)}$ confines $r$
to a single class modulo $2^{\Vi^{(J)}+v+1}$. \qedhere
\end{proof}

% =============================================================================
\section{Density limit and existence criterion}\label{sec:adic}
% =============================================================================

The lift theorem makes the obstruction densities converge to $\cWa$ and
turns the existence question $\cWa > 0$ into a finite diophantine problem;
a closing remark records the $2$-adic Haar-measure reading of $\cWa$.

\begin{proposition}[Density limit]\label{prop:density-limit}
The ratio $\lvert \Obsab \rvert / 2^L$ is non-decreasing in $L$ and bounded
above by $\tfrac12$, so the asymptotic density
\[
\cWab = \lim_{L \to \infty} \frac{\lvert \Obsab \rvert}{2^L}
\]
exists, and every finite ratio $\lvert \Obsab \rvert / 2^L$ is a rigorous
lower bound for it.
\end{proposition}

\begin{proof}
By the lift theorem (Theorem~\ref{thm:lift}),
$h_{L+1} = 2 h_L + \lvert \Atom_{L+1}^{(a,b)} \rvert \ge 2 h_L$ with
$h_L := \lvert \Obsab \rvert$, so $h_{L+1}/2^{L+1} \ge h_L/2^L$. Only odd
residues can be obstructions, so $h_L \le 2^{L-1}$ and the ratio stays
$\le \tfrac12$; a bounded monotone sequence converges. \qedhere
\end{proof}

\begin{definition}[Realisable valuation datum]\label{def:realisable}
A \emph{finite valuation datum} consists of an integer $v \ge 1$, a
termination index $J \ge 0$, and step valuations
$(\vk^{(j)})_{0 \le j < J}$, $(\vi^{(j)})_{0 \le j < J}$, all $\ge 1$. It
is \emph{realisable} if the residue class it pins by the determinacy lemma
(Lemma~\ref{lem:determinacy}) is non-empty and, for some level $L$, its odd
representative $r$ runs the parallel reduction with precisely this valuation
sequence and \emph{terminates at index $J$} (no stop fires before $J$, and a
stop fires at $J$).
\end{definition}

\begin{theorem}[Existence criterion]\label{thm:existence}
The following are equivalent:
\begin{enumerate}[label=\textup{(\arabic*)}]
\item $\cWa > 0$;
\item $\Obsab \ne \emptyset$ for some level $L$;
\item there exists a realisable valuation datum
(Definition~\ref{def:realisable}) satisfying the diophantine
identity~\eqref{eq:diophantine}.
\end{enumerate}
\end{theorem}

\begin{proof}
$(1) \Rightarrow (2)$ is immediate. $(2) \Rightarrow (1)$: if
$r \in \Obs_{L_0}^{(a,b)}$, the lift theorem (Theorem~\ref{thm:lift})
gives $\lvert \Obsab \rvert \ge 2^{L - L_0}$ for every $L \ge L_0$ (both
lifts of an obstruction are obstructions), so $\cWa \ge 2^{-L_0} > 0$ by
Proposition~\ref{prop:density-limit}.

$(2) \Rightarrow (3)$: an obstruction $r \in \Obsab$ runs a finite
reduction terminating at some index $J$ with $\Phi_J = 0$; its valuation
datum is realisable (it is realised by $r$ itself) and
satisfies~\eqref{eq:diophantine} by Theorem~\ref{thm:diophantine}.

$(3) \Rightarrow (2)$: let a realisable datum satisfy~\eqref{eq:diophantine},
and let $r$ be the odd representative pinned by the determinacy lemma
(Lemma~\ref{lem:determinacy}), so that the reduction at $(r, L)$ has exactly
the prescribed valuations and terminates at $J$. Then the recurrence~\eqref{eq:D-recurrence} holds, and the
computation in the proof of Theorem~\ref{thm:diophantine} gives
$D_{J+1} = b\bigl(a^{J+1} - \sum_{j=0}^{J} a^{J-j} X_j\bigr) = 0$ by
\eqref{eq:diophantine}. Since
$\Phi_J(r) \cdot 2^{\Vi^{(J)} + v} = a D_J - b X_J = D_{J+1} = 0$ and
$2^{\Vi^{(J)} + v} \ne 0$, we get $\Phi_J(r) = 0$, so $r \in \Obsab$.
\qedhere
\end{proof}

The reduction $(2) \Leftrightarrow (3)$ recasts existence of obstructions
--- equivalently the conjecture $\cWa > 0$ (Conjecture~\ref{conj:existence})
--- as a concrete diophantine solvability question for the integer
identity~\eqref{eq:diophantine} over realisable valuation data, a question
internal to the finite residue tower rather than one about $\Ztwo$ directly.

\begin{remark}[A $2$-adic Haar reading]\label{rem:haar}
The density has a clean measure-theoretic shadow, in the tradition of the
Bernstein--Lagarias $3x+1$ conjugacy \cite{BernsteinLagarias1996}. Let
$\Ztwo$ be the ring of $2$-adic integers, $\Ztwox$ its units, and $\mu_2$
the normalised Haar measure ($\mu_2(\Ztwo) = 1$, $\mu_2(\Ztwox) = \tfrac12$;
see Robert~\cite{Robert2000}). The cylinders
$S_L := \bigsqcup_{r \in \Obsab}(r + 2^L \Ztwo) \subseteq \Ztwox$ satisfy
$\mu_2(S_L) = \lvert \Obsab \rvert / 2^L$ and increase
($S_L \subseteq S_{L+1}$, since both lifts of an obstruction are
obstructions), so their union
$\Obstad := \bigcup_L S_L = \{\xi \in \Ztwox : \xi \bmod 2^L \in \Obsab
\text{ for all large } L\}$ is open and, by continuity of $\mu_2$ from
below, $\cWab = \mu_2(\Obstad)$. The reduction does \emph{not} extend to a
terminating procedure on $\Ztwox$ directly: at $(a,b) = (3,-1)$ the unit
$\xi = 27$ is an obstruction at every level, yet its $\Qtwo$ K-iterate
$27 \to 5 \to 7 \to \cdots$ never terminates --- so $\Obstad$ is the tail
set above, not the zero-set of an infinite reduction.
\end{remark}

% =============================================================================
\section{\texorpdfstring{$L_{\min}(a)$}{Lmin(a)} and the characterisation of \texorpdfstring{$J = 1$}{J=1} atoms}\label{sec:Lmin}
% =============================================================================

The smallest level at which $\Obsab$ is nonempty is denoted $L_{\min}(a)$
(independent of $b$ by Theorem~\ref{thm:conjugation}). We give a
lower bound and an exact characterisation of the levels at which atoms
with $J = 1$ exist.

\subsection{Lower bound from the order lemma}

\begin{theorem}[Lower bound]\label{thm:Lmin-lower}
For $a$ outside the negative-Mersenne case, $L_{\min}(a) \ge e + 1$.
\end{theorem}

\begin{proof}
By hypothesis $a$ is outside the negative-Mersenne case (whose
$L_{\min}$ is governed by the companion's sub-class analysis,
Section~\ref{sec:neg-mersenne}). The smallest non-empty level is
attained by an atomic obstruction, and
by Theorem~\ref{thm:no-G0-atom-general} (a forward dependency, proved in
Section~\ref{sec:structure}; no result of Section~\ref{sec:structure}
depends on the present bound, so there is no circularity) atomic
obstructions have $\delta_J \ne 0$. By the order lemma
(Lemma~\ref{lem:order-lemma}), $e \mid \delta_J$, hence
$\lvert \delta_J \rvert \ge e$. At termination both cumulative
valuations are bounded by the available modulus,
$\Vk^{(J)} \le L - 1$ and $\Vi^{(J)} + v \le L - 1$, and
$\Vi^{(J)} + v \ge v \ge 1$; therefore, from
$\delta_J = \Vk^{(J)} - (\Vi^{(J)} + v)$,
\[
e \le \lvert \delta_J \rvert \le L - 1.
\]
This gives $L \ge e + 1$, hence $L_{\min}(a) \ge e + 1$. The bound is
not tight. For the Mersenne multipliers $a \in \{3, 7, 15, 31\}$ the
first obstruction sits at $L_{\min}(a) = 2e + 1$ exactly (the $J = 1$
atom of Corollary~\ref{cor:mersenne}); the non-Mersenne multipliers
scatter around this value --- e.g.\ $L_{\min}(5) = 11$ and
$L_{\min}(21) = 10$ against $2e + 1 = 9$ and $13$ respectively --- and the
honest regularity is the asymptotic $L_{\min}(a)/e \to 2$ of
Conjecture~\ref{conj:Lmin-asymptotic}. \qedhere
\end{proof}

\subsection{Characterisation of \texorpdfstring{$J = 1$}{J=1} atoms}

\begin{proposition}[$J = 1$ characterisation]\label{prop:J1-character}
For a positive odd integer $a \ge 3$, the order $e := \ord_{\lvert a \rvert}(2)$ gives
$a \mid 2^e - 1$, so $k_a := (2^e - 1)/a \in \Z$. There exists an atomic
obstruction at some level $L \le 2e + 1$ with termination index $J = 1$ if
and only if the diophantine equation
\begin{equation}\label{eq:J1-criterion}
a - 1 + 2^{v} = 2^{w} k_a
\end{equation}
admits integers $v \ge 1$ and $w$ with $v + 1 \le w \le e$; the atom then
has shift index $s = 1$ and lives at level $L = w + e + 1$, with
$\Vi^{(1)} = w - v$. The principal valuation $v = 1$ specialises
\eqref{eq:J1-criterion} to $a + 1 = 2^{w} k_a$, i.e.\ to
$q_a := (a + 1)/(2 k_a) = 2^{w-1}$ being a power of $2$ with $q_a \ge 2$;
this $v = 1$ solution is the only one for every odd $3 \le a \le 19$. It is
\emph{not} the only one in general: the smallest exception is $a = 21$
($e = 6$, $k_a = 3$), where \eqref{eq:J1-criterion} has no $v = 1$ solution
($q_a = 11/3 \notin 2^{\Z}$) but the $v = 2$ solution $20 + 2^2 = 2^3 \cdot 3$
realises a $J = 1$ atom at level $L = 10$ (the witness residue is
$b$-dependent: $r = 317$ for $b = -1$, with $b = 1$ conjugate
$\Psi_{-1}(317) = 707$). Negative
multipliers are out of scope: the multiplier sign is not neutralisable by
conjugation (cf.\ Appendix~\ref{sec:err-A3}).
\end{proposition}

\begin{proof}
At $J = 1$, with $\Vk^{(0)} = \Vi^{(0)} = 0$ so that $X_0 = 1 - 2^v$, the
diophantine identity~\eqref{eq:diophantine} $a^2 = a X_0 + X_1$ reads
\begin{equation}\label{eq:J1-identity}
a(a - 1) + a\,2^v = 2^{\Vk^{(1)}} - 2^{\Vi^{(1)} + v}.
\end{equation}

\emph{Necessity.} The left-hand side of~\eqref{eq:J1-identity} is positive
(as $a \ge 3$), so $\Vk^{(1)} > \Vi^{(1)} + v$ and the shift
$\delta_1 = \Vk^{(1)} - \Vi^{(1)} - v$ is strictly positive; by the order
lemma (Lemma~\ref{lem:order-lemma}) it is a positive multiple of $e$. For
an atom both cumulative valuations terminate at the boundary of their
modulus, and since $\Vk^{(1)} > \Vi^{(1)} + v$ the binding one is the
K-track, $\Vk^{(1)} = L - 1$, so the hosting level is
$L = \Vk^{(1)} + 1 = \Vi^{(1)} + v + \delta_1 + 1$. With $\Vi^{(1)} \ge 1$
and $v \ge 1$, the target $L \le 2e + 1$ gives
$\delta_1 \le 2e - \Vi^{(1)} - v \le 2e - 2 < 2e$; together with
$e \mid \delta_1$ and $\delta_1 > 0$ this forces $\delta_1 = e$. Writing
$w := \Vi^{(1)} + v$ (so $v + 1 \le w \le e$, the upper bound from
$L \le 2e + 1$) and $\Vk^{(1)} = w + e$, the right-hand side
of~\eqref{eq:J1-identity} collapses to
$2^{w + e} - 2^{w} = 2^{w}(2^e - 1) = 2^{w} k_a\,a$; dividing by $a$ yields
exactly~\eqref{eq:J1-criterion}.

\emph{Sufficiency.} Conversely, suppose $v \ge 1$ and $v + 1 \le w \le e$
solve~\eqref{eq:J1-criterion}; we construct an atom explicitly (the
determinacy lemma gives uniqueness of a realising seed, not its existence,
so we exhibit one). Since $k_a$ is odd and $w > v$, comparing $2$-adic
valuations in $a - 1 + 2^v = 2^w k_a$ forces $\vtwo(a - 1) = v$: a smaller
or larger valuation of $a - 1$ would give $\vtwo(a - 1 + 2^v) < w$. Fix any
odd $b$ and let $r$ be the odd residue determined modulo $2^{w+e+1}$ by
\begin{equation}\label{eq:J1-construct}
a\,r \equiv 2^{w+e} - b \pmod{2^{w+e+1}},
\end{equation}
which is well defined and odd as $a$ is a unit and $2^{w+e} - b$ is odd; by
construction $\vtwo(a r + b) = w + e$. The two remaining valuations follow.
From $a(r + b) = (a r + b) + b(a - 1)$ and $\vtwo(a - 1) = v < w + e$,
\[
\vtwo(r + b) = \vtwo\bigl(a(r + b)\bigr)
            = \vtwo\bigl(2^{w+e} + b(a - 1)\bigr) = v,
\]
so the initial valuation is the prescribed $v$ and $m := (r + b)/2^v$ is
odd. Since $a m + b = \bigl(a(r + b) + b\,2^v\bigr)/2^v$ and, by
\eqref{eq:J1-construct},
$a(r + b) + b\,2^v \equiv 2^{w+e} + b\,(a - 1 + 2^v)
   = 2^{w+e} + b\,2^w k_a \pmod{2^{w+e+1}}$,
\[
a m + b \equiv 2^{w-v}\bigl(2^{e} + b\,k_a\bigr) \pmod{2^{w+e+1-v}},
\qquad \vtwo(a m + b) = w - v,
\]
the cofactor $2^e + b k_a$ being odd. Hence the parallel reduction of $r$ at
level $L := w + e + 1 \le 2e + 1$ takes one step with $\vk^{(0)} = w + e$
and $\vi^{(0)} = w - v$, reaching $\Vk^{(1)} = w + e = L - 1$,
$\Vi^{(1)} = w - v$; the residual K-modulus then has valuation
$L - \Vk^{(1)} = 1$, so the K-track stops and $J = 1$ (the I-step is
admissible because $w - v < w + e + 1 - v$). These valuations
satisfy~\eqref{eq:J1-identity} ($\delta_1 = e$), so the recurrence
\eqref{eq:D-recurrence} gives $D_{J+1} = 0$ exactly as in the proof of the
existence criterion (Theorem~\ref{thm:existence}), i.e.\ $\Phi_1 = 0$ and
$r \in \Obsab$. Finally $\Vk^{(1)} = L - 1$ is the atom boundary
(Proposition~\ref{prop:atom-J}, Step~2): at level $L - 1$ the K-track stops
already at index $0$, so $r \bmod 2^{L-1} \notin \Obs_{L-1}^{(a,b)}$ and $r$
is atomic. The smallest non-$v{=}1$ instance is $a = 21$ ($v = 2$, $w = 3$),
giving the $b = -1$ witness $r = 317$ at $L = 10$ ($b = 1$ conjugate $707$;
Table~\ref{tab:J1}).
\qedhere
\end{proof}

\begin{corollary}[Mersenne construction]\label{cor:mersenne}
Every positive Mersenne multiplier $a = 2^e - 1$ admits $J = 1$ atoms at
levels $L \le 2e + 1$. (The negative Mersenne multipliers $a = -(2^e - 1)$
also carry $J = 1$ atoms, but through the distinct sub-class mechanism of
the companion manuscript --- where they arise alongside the $J = 0$
obstructions absent for positive $a$ --- and are not instances of this
corollary; see Section~\ref{sec:neg-mersenne}.)
\end{corollary}

\begin{proof}
For $a = 2^e - 1$, $k_a = 1$ and $q_a = (a + 1)/2 = 2^{e-1}$, a power of
$2$ of value $\ge 2$ when $e \ge 2$. Apply
Proposition~\ref{prop:J1-character}. \qedhere
\end{proof}

The solutions of the criterion~\eqref{eq:J1-criterion} and the smallest
level realising $J = 1$ for small odd $\lvert a \rvert$ are collected in
Table~\ref{tab:J1}.

\begin{table}[ht]
\centering
\begin{tabular}{rrrll}
\toprule
$\lvert a \rvert$ & $e$ & $k_a$ & $(v, w)$ solving \eqref{eq:J1-criterion} & $L^{J=1} = w + e + 1$ \\
\midrule
3 & 2 & 1 & $(1, 2)$ & 5 \\
5 & 4 & 3 & none & --- ($\ge 11$ empirical) \\
7 & 3 & 1 & $(1, 3)$ & 7 \\
9 & 6 & 7 & none & --- ($\ge 11$ empirical) \\
11 & 10 & 93 & none & --- ($\ge 18$ empirical) \\
15 & 4 & 1 & $(1, 4)$ & 9 \\
17 & 8 & 15 & none & --- ($\ge 16$ empirical) \\
21 & 6 & 3 & $(2, 3)$ & 10 \\
31 & 5 & 1 & $(1, 5)$ & 11 \\
\bottomrule
\end{tabular}
\caption{Solvability of the $J = 1$ criterion~\eqref{eq:J1-criterion}
$a - 1 + 2^v = 2^w k_a$ for small odd $\lvert a \rvert$. Entries marked
``none'' have no admissible $(v, w)$, so $J \ge 2$ throughout (their
empirical $L_{\min}$ from direct enumeration up to $L = 18$ is shown). The
Mersenne multipliers $3, 7, 15, 31$ realise $v = 1$; $a = 21$ is the
smallest multiplier whose only solution has $v \ge 2$
(\srcpath{code/python/count_obstructions_ac.py}).}
\label{tab:J1}
\end{table}

\begin{conjecture}[Asymptotic $L_{\min}$]\label{conj:Lmin-asymptotic}
$L_{\min}(a) \le 2e + O(1)$ as $\lvert a \rvert \to \infty$, and
$L_{\min}(a)/e \to 2$.
\end{conjecture}

The conjecture is consistent with all tested $\lvert a \rvert \in \{3,
\ldots, 31\}$; a rigorous proof requires either a characterisation of
$J = j$ solvability for $j \ge 2$ or a Hensel-type lifting argument in
$\Ztwo$.

\subsection{Existence for non-Mersenne multipliers}

\begin{conjecture}[Existence]\label{conj:existence}
$\cWa > 0$ for every odd $a$ with $\lvert a \rvert \ge 3$.
\end{conjecture}

By Theorem~\ref{thm:existence}, this reduces to realisability of the
diophantine identity~\eqref{eq:diophantine} for at least one tuple.
Empirically, $\cWa > 0$ is confirmed for $\lvert a \rvert \in \{3, 5,
7, 9, 15, 17, 21, 31\}$. For $\lvert a \rvert \in \{11, 13, 19, 23, 25,
27\}$ no obstruction is detected at levels $L \le 16$, but a deeper
enumeration locates the first obstruction at progressively higher
levels --- $L = 18$ for $\lvert a \rvert = 11$, $L = 24$ for
$\lvert a \rvert = 13$, and $L = 26$ for $\lvert a \rvert = 23$ ---
while $\lvert a \rvert \in \{19, 25, 27\}$ remain undetected through
$L = 26$; this growing ``waiting time'' is consistent with a longer
onset rather than vacuity. The enumeration
underlying these statements (and the tables of
Sections~\ref{sec:Lmin}--\ref{sec:scaling}) is carried out for
\emph{positive} $a$ together with the degenerate $a = -1$
(Section~\ref{sec:edge}); by the conjugation theorem the cardinalities
are $b$-independent, but the negative non-degenerate multipliers
$a \le -3$ are not separately enumerated here --- the negative-Mersenne
sub-family is treated in the companion manuscript
(Section~\ref{sec:neg-mersenne}), and the structural results of
Section~\ref{sec:structure} cover negative $a$ through their proofs (with
the negative-Mersenne caveat noted in the proof of
Theorem~\ref{thm:no-G0-atom-general}).

% =============================================================================
\section{Structural results over \texorpdfstring{$a$}{a}}\label{sec:structure}
% =============================================================================

\subsection{Generalised no-shift-zero-atom theorem}

\begin{theorem}[No shift-zero atom, generalised]\label{thm:no-G0-atom-general}
Let $a$ be an odd integer with $\lvert a \rvert \ge 3$ that is \emph{not}
a negative Mersenne number $-(2^{v_M} - 1)$ ($v_M \ge 2$), let $b$ be an
odd nonzero integer, and $e := \ord_{\lvert a \rvert}(2)$. For every level
$L$,
\[
\lvert \Atom_L^{(a,b), G_0} \rvert = 0;
\]
equivalently, every atomic obstruction has nonzero shift index.
\end{theorem}

\begin{proof}
The strategy, adapting \cite[Theorem~6.2]{ObstructionResidues2026}, is to
show that \emph{every} shift-zero obstruction is a lift of a level-$(L-1)$
obstruction, hence not atomic. The prior argument carries over under
$3 \to a$, $-1 \to b$ with the parity step $\delta \in 2\Z$ replaced by
$\delta \in e\Z$ (order $e := \ord_{\lvert a \rvert}(2)$). The prior argument
already forces $J \ge 2$, but trivially --- for $a = 3$ no shift-zero
certificate can vanish at $J \le 1$ for any $v$; what genuinely changes here
is the \emph{content} of that step: the general equation $2^v = 1 - a$ is
solvable precisely for the negative-Mersenne multipliers, which it thereby
isolates and forces us to exclude.

\emph{$J \ge 2$ is forced (a shift-zero obstruction needs at least two
steps).} A shift-zero obstruction has $\delta_J = 0$. At $J = 0$ this is
impossible: $\delta_0 = \Vk^{(0)} - \Vi^{(0)} - v = -v \le -1 \ne 0$
(note this is the correct obstruction to $J = 0$ --- the certificate value
$\Phi_0$ itself can vanish for some negative $a$, but never with
$\delta_0 = 0$). At $J = 1$ with $\delta_1 = 0$ we have $c_1 = 1$, so the
certificate $\Phi_1 = a d_1$ vanishes only if $d_1 = 0$; the $d$-update
then forces $a d_0 + b(1 - c_0) = 0$, i.e.\ $(a - 1)\,2^{-v} = -1$, i.e.\
$2^{v} = 1 - a$. For $a > 0$ the right-hand side is $\le -2$, impossible.
For $a < 0$ it requires $\lvert a \rvert + 1 = 2^{v}$, that is $\lvert a
\rvert$ a Mersenne number; the negative-Mersenne multipliers are treated
by the finer sub-class analysis of the companion manuscript
(Section~\ref{sec:neg-mersenne}), while for negative non-Mersenne $a$ the
equation $2^{v} = \lvert a \rvert + 1$ has no solution. Hence, outside the
negative-Mersenne case, no shift-zero obstruction has $J \le 1$, so
$J \ge 2$.

\emph{The projected reduction.} Let $r \in \Obs_L^{(a,b), G_0}$ and set
$r' := r \bmod 2^{L-1}$. The parallel reductions at $(r, L)$ and
$(r', L-1)$ start from the same residue classes with the modulus halved;
exactly as in the lift theorem (Theorem~\ref{thm:lift}) the two
K-tracks differ by $\Delta_K^{(j)} = 2^{L-1-\Vk^{(j)}} a^j$ (and the
I-tracks by $\Delta_I^{(j)} = 2^{L-1-v-\Vi^{(j)}} a^j$; $a$ odd keeps the
$2$-adic valuation). While $\Vk^{(j)} < L-1$ the ultrametric identity
forces the per-step valuations, hence the affine data $(c_j, d_j)$, to
agree between the two reductions. Since $J \ge 2$ the case split below is
well defined.

\emph{Case A ($\Vk^{(J)} < L-1$).} No level-$(L-1)$ stop fires before
index $J$: a K-stop at index $J' \le J-1$ would force
$\Vk^{(J)} \ge \Vk^{(J'+1)} \ge L-1$, and an I-stop at index $J' \le J-1$
would force $\Vi^{(J)} \ge \Vi^{(J'+1)} \ge L-1-v$, which under
$\delta_J = 0$ (so $\Vi^{(J)} = \Vk^{(J)} - v$) again reads
$\Vk^{(J)} \ge L-1$ --- both contradicting the Case-A hypothesis. Hence the
level-$(L-1)$ reduction reaches index $J$ with the same terminal data
$(c_J, d_J)$, so $\Phi_J(r') = 0$ and $r' \in \Obs_{L-1}^{(a,b)}$ (in the
same shift class $G_0$).

\emph{Case B ($\Vk^{(J)} = L-1$).} Now the level-$(L-1)$ K-stop fires one
step earlier, at index $J-1$ --- and not before: for any $J' \le J-2$ one
has $\Vk^{(J'+1)} \le \Vk^{(J-1)} < \Vk^{(J)} = L-1$, and under $\delta_J = 0$
(so $\Vi^{(J)} = \Vk^{(J)} - v$) likewise
$\Vi^{(J'+1)} \le \Vi^{(J-1)} < L-1-v$, so no level-$(L-1)$ stop precedes
index $J-1$. The $(r', L-1)$ reduction therefore terminates at index $J-1$
with data $(c_{J-1}, d_{J-1})$. From the $d$-update (Lemma~\ref{lem:affine}) the
parent's certificate value is
$d_J = \bigl(a\,d_{J-1} + b(1 - c_{J-1})\bigr)/2^{\vi^{(J-1)}}$, and the
$G_0$ condition $d_J = 0$ gives $a\,d_{J-1} + b(1 - c_{J-1}) = 0$, i.e.\
\[
\Phi_{J-1}(r') = b(1 - c_{J-1}) + a\,d_{J-1} = 0.
\]
Thus $r'$ terminates as an obstruction at level $L-1$, in shift class
$\delta_{J-1}(r')/e$ (an integer by the order lemma,
Lemma~\ref{lem:order-lemma}), so $r' \in \Obs_{L-1}^{(a,b)}$.

In both cases $r' \in \Obs_{L-1}^{(a,b)}$, so $r$ --- being one of the two
lifts of $r'$ --- is not atomic (Theorem~\ref{thm:lift}). Hence no
shift-zero obstruction is atomic, $\lvert \Atom_L^{(a,b), G_0} \rvert = 0$.
The $J \ge 2$ forcing above is purely algebraic and holds at \emph{every}
level; the case split then applies at every $L$ for which the projected
level-$(L-1)$ reduction is defined, so the conclusion holds at all such
$L$. The finitely many smallest levels are confirmed by direct
enumeration, which gives $\lvert \Atom_L^{(a,b), G_0} \rvert = 0$ for all
odd $3 \le a \le 31$ (and the degenerate $a = -1$) up to $L = 16$; the
negative non-degenerate multipliers are covered by the argument, not the
enumeration. \qedhere
\end{proof}

\begin{corollary}[Synchronisation, generalised]\label{cor:sync-general}
For every $r \in \Obsab$ both tracks terminate at the same index $J$, and
the cumulative valuations are synchronised to an integer multiple of $e$:
\[
\Vk^{(J)} - \bigl(\Vi^{(J)} + v\bigr) = \delta_J = e \cdot s, \qquad s := \delta_J / e \in \Z.
\]
The $\Tneg$-manuscript's case is $a = 3$, $e = 2$.
\end{corollary}

\begin{proof}
The common termination index and $\vi^{(J)} = \vk^{(J)} + \delta_J$ are the
terminal synchronisation (Remark~\ref{rem:sync}); the displayed equality is
the definition $\delta_J = \Vk^{(J)} - \Vi^{(J)} - v$, and $e \mid \delta_J$
is the order lemma (Lemma~\ref{lem:order-lemma}). \qedhere
\end{proof}

\subsection{Polynomial atom bound at fixed \texorpdfstring{$J$}{J}}

\begin{proposition}[Atom bound at fixed $J$]\label{prop:atom-J}
For odd $a$ with $\lvert a \rvert \ge 3$, odd $b \ne 0$, and every
$J \ge 1$, the number of atomic $G_{\ne 0}$-obstructions at level $L$ with
termination index $J$ satisfies
\[
\lvert \Atom_L^{(a,b), G_{\ne 0}} \rvert_J
   \;\le\; \binom{L-2}{J-1} + \binom{L-2}{J}
   \;=\; \binom{L-1}{J}
   \;=\; O_J\!\bigl(L^{J}\bigr),
\]
where $\lvert \cdot \rvert_J$ counts atomic $G_{\ne 0}$-obstructions with
the given $J$. In particular every fixed-$J$ stratum is polynomially
bounded in $L$. The bound is unconditional; no per-step factor is left open.
\end{proposition}

\begin{proof}
The proof has three steps: a $2$-adic determinacy lemma (a fixed valuation
sequence pins down the residue), the atom boundary condition, and a count.

\emph{Step 1 (determinacy).} By the determinacy lemma
(Lemma~\ref{lem:determinacy}), fixing the K-sequence
$\vk^{(0)}, \ldots, \vk^{(J-1)}$ pins $r$ to at most one residue class
modulo $2^{\Vk^{(J)}+1}$, and fixing $v$ together with
$\vi^{(0)}, \ldots, \vi^{(J-1)}$ pins it to at most one class modulo
$2^{\Vi^{(J)}+v+1}$.

\emph{Step 2 (atom boundary).} At termination both cumulative valuations
obey $\Vk^{(J)} \le L-1$ and $\Vi^{(J)} + v \le L-1$ (a successful step $j$
requires $\vk^{(j)} < L - \Vk^{(j)}$, resp.\ $\vi^{(j)} < L - v - \Vi^{(j)}$,
so $\Vk^{(j+1)} \le L-1$, resp.\ $\Vi^{(j+1)} + v \le L-1$). Suppose both
were strict: $\Vk^{(J)} < L-1$ and $\Vi^{(J)} + v < L-1$. Then the
discrepancy argument from the proof of the lift theorem
(Theorem~\ref{thm:lift}) applies to $r' := r \bmod 2^{L-1}$: the two
reductions differ by $\Delta_K^{(j)} = 2^{L-1-\Vk^{(j)}} a^j$ (and the
I-analogue), and because $a$ is odd and $\Vk^{(j)} \le \Vk^{(J)} < L-1$
(resp.\ $\Vi^{(j)} + v \le \Vi^{(J)} + v < L-1$) throughout, the ultrametric
identity forces equal per-step valuations and hence equal terminal data, so
$\Phi_J(r') = 0$ and $r' \in \Obs_{L-1}^{(a,b)}$. Then $r$, a lift of $r'$,
is not atomic. Contrapositively, every atomic $G_{\ne 0}$-obstruction
satisfies the boundary condition
\[
\Vk^{(J)} = L-1 \qquad\text{or}\qquad \Vi^{(J)} + v = L-1.
\]

\emph{Step 3 (count).} Combine Steps~1 and~2. If $\Vk^{(J)} = L-1$, the
K-sequence pins $r$ modulo $2^{\Vk^{(J)}+1} = 2^{L}$, hence to at most one
residue in $\{1, \ldots, 2^L\}$; the admissible K-sequences are the
compositions of $L-1$ into $J$ positive parts, $\binom{L-2}{J-1}$ in number.
If instead $\Vi^{(J)} + v = L-1$, the datum $(\vi^{(0)}, \ldots,
\vi^{(J-1)}, v)$ pins $r$ modulo $2^{\Vi^{(J)}+v+1} = 2^{L}$, again to at
most one residue; these data are the compositions of $L-1$ into $J+1$
positive parts, $\binom{L-2}{J}$ in number. By Step~2 every atom lies in at
least one case, and by Step~1 each datum contributes at most one residue, so
\[
\lvert \Atom_L^{(a,b), G_{\ne 0}} \rvert_J
  \;\le\; \binom{L-2}{J-1} + \binom{L-2}{J} \;=\; \binom{L-1}{J}
\]
by Pascal's rule. Since $\binom{L-1}{J} \le (L-1)^J/J!$, this is
$O_J(L^J)$. \qedhere
\end{proof}

\begin{remark}
The exponent $J$ is generous: the boundary congruences couple the two
tracks, and the order lemma (Lemma~\ref{lem:order-lemma}) and the
realisability of each composition cut the admissible data far below
$\binom{L-1}{J}$ (e.g.\ at $a = 3$ the $J = 5$ stratum peaks near
$30$ rather than $\binom{L-1}{5}$). What the bound establishes
unconditionally is that each fixed-$J$ stratum is polynomial in $L$; the
\emph{summation} over $J$ --- where the strata overlap and $J$ itself
ranges up to $J_{\max} = L - 4$ --- is the open content of the
growth-exponent question of Section~\ref{sec:spectral}
(Conjecture~\ref{conj:alpha}, and the open problem whether $\alpha(a) = 1$).
\end{remark}

Summed over all $J$, the total atom count $\lvert \Atom_L^{(a,b),
G_{\ne 0}} \rvert$ fixes the growth exponent
$\alpha(a) := \limsup_L \lvert \Atom_L \rvert^{1/L} \in [1, 2]$. Whether
this total is genuinely polynomial --- the \emph{strong} reading
$\alpha(a) = 1$ --- or only sub-$2$ exponential ($\alpha(a) \in (1, 2)$)
we leave as an \emph{open problem} rather than a conjecture: the deeper
empirical data of Section~\ref{sec:spectral} lean against the polynomial
reading. What we do conjecture is the strictly weaker bound
$\alpha(a) < 2$ (Conjecture~\ref{conj:alpha} in
Section~\ref{sec:spectral}), which already suffices for the series
representation of Corollary~\ref{cor:series} to converge and is the
principal \emph{tractable} target.

Proposition~\ref{prop:atom-J} controls each fixed-$J$ stratum
polynomially, but the summation over $J$ --- where $J$ ranges up to
$J_{\max} = L - 4$ --- would need a concentration estimate on the
$J$-distribution that we do not have. The early hypothesis
$J_{\max}(L) = O(\log L)$ is refuted by Theorem~\ref{thm:Lminus4-atom}
below (Appendix~\ref{sec:err-Jmax}); the corrected picture is a bell-shaped
$J$-distribution with mode
$\approx \log_2 L + e$ and a heavy tail to $J_{\max} = L - 4$
($\lvert a \rvert = 3$). For $a = 3$ the polynomial question is identical
to the $\Tneg$-manuscript's open problem
(\cite[Section~6.4]{ObstructionResidues2026}), here extended to all
$\lvert a \rvert \ge 3$.

\subsection{Constructive extremal family at \texorpdfstring{$J = L - 4$}{J=L-4} for \texorpdfstring{$a = 3$}{a=3}}

\begin{theorem}[$L - 4$ extremal family]\label{thm:Lminus4-atom}
For $a = 3$, $b = -1$, and every \emph{even} $L \ge 6$, the residue
$r := 2^{L-2} + 3$ is an atomic obstruction in $\Obs_L^{(3,-1)}$ with
termination index $J = L - 4$ and shift index $s = 1$. (For odd $L$ this
particular residue is not an obstruction --- the trace below yields
$\delta_J = 1$, an odd shift index forbidden by the order lemma; the
extremal index for odd $L$ is not addressed here.)
\end{theorem}

\begin{proof}
We exhibit the iteration trace. At $j = 0$: $r = 2^{L-2} + 3$,
$3r - 1 = 3 \cdot 2^{L-2} + 8 = 2^3 (3 \cdot 2^{L-5} + 1)$, so
$\vk^{(0)} = 3$; $m = (r - 1)/2 = 2^{L-3} + 1$ (with $v = 1$),
$3m - 1 = 3 \cdot 2^{L-3} + 2 = 2(3 \cdot 2^{L-4} + 1)$, so $\vi^{(0)} =
1$.

At $j = 1, \ldots, L - 6$: $a_K^{(j)} = 3^j \cdot 2^{L-4-j} + 1$,
$3 a_K^{(j)} - 1 = 3^{j+1} \cdot 2^{L-4-j} + 2 = 2(3^{j+1} \cdot 2^{L-5-j} + 1)$,
so $\vk^{(j)} = 1$; analogously $\vi^{(j)} = 1$.

At $j = L - 5$: $a_K^{(L-5)} = 2 \cdot 3^{L-5} + 1$, $3 \cdot
a_K^{(L-5)} - 1 = 2 (3^{L-4} + 1)$. We take $L$ even, so $3^{L-4} \equiv
1 \pmod 4$ and $\vk^{(L-5)} = 2$. (For odd $L$ one has $\vk^{(L-5)} = 1$
instead; the cumulative valuation then yields $\delta_J = 1$, an odd
shift index, so $r$ fails the order lemma and is not an obstruction ---
hence the even-$L$ restriction in the statement.)

At $j = L - 4$: $\Vk^{(L-4)} = 3 + (L - 6) \cdot 1 + 2 = L - 1$, so
$\vtwo(b_K^{(L-4)}) = L - \Vk^{(L-4)} = 1$, and the K-track terminates
at this step (assuming $\vk^{(L-4)} \ge 1$, which holds). On the I-track,
$\Vi^{(L-4)} = (L - 4) \cdot 1 = L - 4$, and the residual modulus
$\vtwo(b_I^{(L-4)}) = L - v - \Vi^{(L-4)} = L - 1 - (L-4) = 3$, so the
I-track also terminates. Both at $J = L - 4$.

Shift index: $\delta_J = \Vk^{(J)} - \Vi^{(J)} - v = (L - 1) - (L - 4) - 1
= 2 = 1 \cdot e$, so $s = 1$. Certificate: $\Phi_J = b(1 - c_J) + a d_J$
with $c_J = 2^2 = 4$ and $d_J$ computable from the diophantine identity
to $(1 - c_J)/3 = -1$; $\Phi_J = (-1)(1 - 4) + 3 \cdot (-1) = 3 - 3 = 0$.
So $r \in \Obs_L^{(3, -1)}$.

Atomicity: run the reduction on $r \bmod 2^{L-1}$ at level $L - 1$. It
agrees with the level-$L$ trace until the K-track reaches
$\Vk^{(L-5)} = 3 + (L-6) = L - 3$, where the residual modulus is
$(L-1) - (L-3) = 2$; since $\vk^{(L-5)} = 2 \not< 2$, the K-track stops one
step early, at index $L - 5$. There the shift index is
\[
\delta_{L-5} = \Vk^{(L-5)} - \Vi^{(L-5)} - v = (L-3) - (L-5) - 1 = 1,
\]
an \emph{odd} value (here $\Vi^{(L-5)} = L - 5$ from $\vi^{(j)} = 1$,
$j \le L - 6$). By the order lemma (Lemma~\ref{lem:order-lemma}, $e = 2$)
no obstruction can have odd shift, so $\Phi_{L-5}(r \bmod 2^{L-1}) \ne 0$
and $r \bmod 2^{L-1} \notin \Obs_{L-1}^{(3,-1)}$. Hence $r$ is atomic.
\qedhere
\end{proof}

\begin{conjecture}[Maximality of $L - 4$ for $a = 3$]\label{conj:Jmax}
For $a = 3$, $b = -1$, and every \emph{even} $L \ge 6$, $J = L - 4$ is
the maximum termination index attained by any atomic obstruction in
$\Obs_L^{(3,-1)}$.
\end{conjecture}

Empirically verified for even $L \in \{10, 12, 14, 16\}$
(\srcpath{code/python/count_obstructions_ac.py}). A rigorous proof
would require non-existence statements at $J = L - 3, L - 2, L - 1$,
which appear to follow from the diophantine identity reduced modulo
$\lvert a \rvert^{J+1}$ but we do not pursue this here.

\subsection{\texorpdfstring{$J$}{J}-distribution: empirical picture}

Empirical $J$-statistics at $a = 3$, $b = -1$:

\begin{table}[ht]
\centering
\begin{tabular}{rrrrr}
\toprule
$L$ & $\lvert \Atom_L \rvert$ & $J_{\min}$ & $J_{\max}$ & $\bar J$ \\
\midrule
10 & 7 & 3 & 6 & 4.00 \\
12 & 21 & 4 & 8 & 4.76 \\
14 & 66 & 4 & 10 & 5.64 \\
16 & 214 & 4 & 12 & 6.57 \\
18 & 722 & 5 & 14 & --- \\
\bottomrule
\end{tabular}
\caption{Termination-index statistics for atomic
$G_{\ne 0}$-ob\-struc\-tions at $a = 3$, $b = -1$. $J_{\max} = L - 4$ matches
Theorem~\ref{thm:Lminus4-atom}; the mean $\bar J$ grows
slowly (logarithmically) in $L$ --- broadly consistent with a $\log_2 L +
e$ trend, though the few available points do not pin the constant down.}
\label{tab:J-stats}
\end{table}

Across $\lvert a \rvert \in \{3, 5, 7\}$ at $L = 16$, the gap
$L - J_{\max}(L, a)$ between $L$ and the largest termination index of an
\emph{atomic} $G_{\ne 0}$-obstruction (as in Table~\ref{tab:J-stats})
depends on $a$ (Table~\ref{tab:Jmax-a}):

\begin{table}[ht]
\centering
\begin{tabular}{rrrrr}
\toprule
$\lvert a \rvert$ & $e$ & $L$ & $J_{\max}$ & $L - J_{\max}$ \\
\midrule
3 & 2 & 16 & 12 & 4 \\
5 & 4 & 16 & 6 & 10 \\
7 & 3 & 16 & 5 & 11 \\
\bottomrule
\end{tabular}
\caption{Atomic $J_{\max}$ across $\lvert a \rvert$ at $L = 16$, taken --- as
in Table~\ref{tab:J-stats} --- over atomic $G_{\ne 0}$-obstructions (the
maximum over \emph{all} obstructions is larger, e.g.\ $8$ for $a = 5$ and
$10$ for $a = 7$). The $\Tneg$-manuscript value $L - J_{\max} = 4 = 2e$ at
$a = 3$ does not generalise as $2e$; see Conjecture~\ref{conj:Jmax-scaling}.}
\label{tab:Jmax-a}
\end{table}

\begin{conjecture}[$L - J_{\max}$ scaling]\label{conj:Jmax-scaling}
$L - J_{\max}(L, a) \approx \mathrm{const}(a) \cdot e + O(1)$ with
$\mathrm{const}(3) = 2$.
\end{conjecture}

% =============================================================================
\section{The growth exponent \texorpdfstring{$\alpha(a)$}{alpha(a)}}\label{sec:spectral}
% =============================================================================

\subsection{\texorpdfstring{$\alpha(a) < 2$}{alpha(a) < 2}}

Let $\alpha(a) := \limsup_L \lvert \Atom_L^{(a), G_{\ne 0}} \rvert^{1/L}$.
The series convergence in Corollary~\ref{cor:series} gives $\alpha(a)
\le 2$ trivially.

\begin{conjecture}[$\alpha(a) < 2$ strict]\label{conj:alpha}
For every odd $a$ with $\lvert a \rvert \ge 3$, $\alpha(a) < 2$.
\end{conjecture}

This is the weak form of the growth-exponent question
($\alpha(a) = 1 \Rightarrow \alpha(a) < 2$) and is equivalent to
$\lvert \Atom_L \rvert / 2^L \to 0$. Direct enumeration constrains
$\alpha(3)$ without determining it: over $a = 3$ (Table~\ref{tab:J-stats}
with $\lvert \Atom_L \rvert = 7, 21, 66, 214, 722$ at
$L = 10, 12, \ldots, 18$, extended to $\lvert \Atom_{24} \rvert = 32051$),
the raw root $\lvert \Atom_L \rvert^{1/L}$ rises $1.44 \to 1.54$
($L = 18, 24$) and the per-step ratio
$\lvert \Atom_L \rvert / \lvert \Atom_{L-1} \rvert$ rises
$1.85 \to 1.91$, with no sign of the turnover that the polynomial regime
$\alpha(3) = 1$ would eventually require (a polynomial $L^d$ has its ratio
\emph{decreasing} to $1$). The data are therefore more readily read as
$\alpha(3) \in (1, 2)$ than as $\alpha(3) = 1$, though a finite
enumeration cannot decide a $\limsup$; this is why we state only
$\alpha(a) < 2$ as a conjecture and leave $\alpha(a) = 1$ as an open
problem.

\subsection{A possible spectral route}

One natural attack on Conjecture~\ref{conj:alpha} adapts Wirsching's
Markov operator for $\Tpos = (3n+1)/2^v$ (\cite{Wirsching1998}), an
$L^1(\Ztwo)$-contraction with a unique invariant density. The analogue
here would act on the parallel-reduction state $(c, d)$ through the update
rule of Lemma~\ref{lem:affine}, weighting the transition with step
valuations $(v_K, v_I)$ by $2^{-v_K - v_I}$; a spectral gap $\lambda_1 < 1$
of such an operator would give $\alpha(a) = 2\lambda_1 < 2$. We record this
only as a possible route and do not develop it: unlike Wirsching's setting,
the state space $(c, d) \in \Qtwo \times \Qtwo$ is unbounded, and neither
the existence of an invariant density nor a spectral gap is clear.

% =============================================================================
\section{Empirical density scaling}\label{sec:scaling}
% =============================================================================

We report the empirical scaling of $\cWa$ with $\lvert a \rvert$ via
direct enumeration at level $L \le 16$
(\srcpath{code/python/count_obstructions_ac.py}). The density estimates
in Table~\ref{tab:scaling} are restricted to the multipliers
$\lvert a \rvert \in \{3, 5, 7, 15, 31\}$ that are well resolved at
$L = 16$. The remaining tested multipliers do carry obstructions at
$L = 16$ ($\lvert\Obs_{16}^{(9,-1)}\rvert = 47$,
$\lvert\Obs_{16}^{(21,-1)}\rvert = 69$,
$\lvert\Obs_{16}^{(17,-1)}\rvert = 1$), confirming $\cWa > 0$, but with
too few obstructions for a stable density estimate at this level
(for $\lvert a \rvert = 17$ a single obstruction); they are therefore
omitted from the scaling fit.

\begin{table}[ht]
\centering
\begin{tabular}{rrrrr}
\toprule
$\lvert a \rvert$ & $e$ & $4^{-e}$ & $\cWa$ (empir.\ lower bound, $L = 16$) & $C(a) := \cWa / 4^{-e}$ \\
\midrule
3 & 2 & $0.0625$ & $0.1153$ & $1.84$ \\
7 & 3 & $0.0156$ & $0.0146$ & $0.93$ \\
15 & 4 & $0.00391$ & $0.00243$ & $0.62$ \\
5 & 4 & $0.00391$ & $0.00172$ & $0.44$ \\
31 & 5 & $0.000977$ & $0.00053$ & $0.55$ \\
\bottomrule
\end{tabular}
\caption{Empirical density scaling over the five multipliers well
resolved at $L = 16$. Each $\cWa$ entry is a rigorous lower bound for the
limit (the sequence $\lvert \Obsab \rvert / 2^L$ is non-decreasing, by the
lift theorem; see Proposition~\ref{prop:density-limit}). A log-linear least-squares fit
gives $\log_2 \cWa \approx -2.64\,e + 1.91$ ($R^2 = 0.99$). The fitted
slope $-2.64$ is somewhat steeper than the $-2$ that exact $4^{-e}$ scaling
would produce: over these five points the data are consistent both with a
bounded correction factor $C(a) := \cWa / 4^{-e}$ (which stays in $[0.44,
1.84]$ here) and with a slowly \emph{decaying} $C(a)$, i.e.\ a decay faster
than $4^{-e}$. We therefore report $\cWa \approx C(a)\,4^{-e}$ as an
empirical law and leave the asymptotic boundedness of $C(a)$ --- the
two-sided form $\cWa = \Theta(4^{-e})$ --- as Conjecture~\ref{conj:correction}.}
\label{tab:scaling}
\end{table}

The structural mechanism behind a $4^{-e}$ \emph{lower} bound is the lift
theorem (Theorem~\ref{thm:lift}): a single atomic obstruction at
level $L_{\min}$ produces $2^{L - L_{\min}}$ obstructions at every level
$L \ge L_{\min}$, so $\cWa \gtrsim 2^{-L_{\min}(a)}$. Granting the
(conjectural) $L_{\min}(a) \approx 2e$ of
Conjecture~\ref{conj:Lmin-asymptotic}, this gives
\[
\cWa \gtrsim 2^{-L_{\min}(a)} \approx 4^{-e}.
\]
Note that this argument bounds $\cWa$ only from below; the matching upper
bound is the open content of Conjecture~\ref{conj:correction}.

\begin{conjecture}[Bounded correction factor]\label{conj:correction}
For \emph{positive} odd $a \ge 3$ there exist constants $C_0 > 0$,
$C_1 < \infty$ with $C_0 \le C(a) \le C_1$.
\end{conjecture}

The restriction to positive $a$ is essential, not cosmetic. For the
negative Mersenne multipliers $a = -(2^{v_M} - 1)$ the density scales as
$\cWa = \Theta(2^{-e})$ (Section~\ref{sec:neg-mersenne}), the square root
of the positive scale, so $C(a) = \cWa/4^{-e} = \Theta(2^{e}) \to \infty$;
the bound therefore cannot hold across the sign axis. Already the smallest
negative non-Mersenne multiplier $a = -5$ ($e = 4$) leaves the
positive-$a$ band: the $b$-independent cardinality
$\lvert \Obs_{14}^{(-5,b)} \rvert = 357$ gives the rigorous lower bound
$C(-5) \ge 357/64 > 5.5$, against the $[0.44, 1.84]$ of
Table~\ref{tab:scaling}. The empirical law $\cWa \approx C(a)\,4^{-e}$ and
its two-sided form $\cWa = \Theta(4^{-e})$ are thus phenomena of the
\emph{positive} multipliers; the negative multipliers are governed by the
companion's sub-class analysis (Section~\ref{sec:neg-mersenne}) for the
Mersenne case and are left uncharacterised otherwise.

\begin{remark}[Deeper enumeration and extrapolation]\label{rem:deep-scaling}
The $C(a)$ entries of Table~\ref{tab:scaling} are rigorous lower bounds at
the truncation level $L = 16$; they understate the limit, and the gap is
already visible a few levels deeper. Pushing the (unhacked) enumeration to
$L = 24$ --- still a rigorous lower bound by
Proposition~\ref{prop:density-limit} --- raises $C(3)$ from $1.84$ to
$2.15$, and over the eight multipliers
$\lvert a \rvert \in \{3, 5, 7, 9, 15, 17, 21, 31\}$ the lower-bound band
widens to $[0.56, 4.40]$ (extremes at $a = 31$ and $a = 21$), already
beyond the $[0.44, 1.84]$ read off at $L = 16$. Fitting the atomic series
$\lvert \Atom_l^{(a), G_{\ne 0}} \rvert \approx K_a\,\alpha(a)^l$ on the
deepest available window and summing the geometric tail
$\sum_{l > 24} \lvert \Atom_l \rvert / 2^l$ yields \emph{heuristic} limit
estimates that stay in a comparable bounded band (e.g.\ $C(3) \approx 2.6$,
$C(21) \approx 4.4$), with fitted growth $\alpha(a) < 2$ throughout
($\alpha(a) \in [1.46, 1.88]$ for the well-resolved $a \le 15$; the sparser
high-order multipliers are consistent but under-resolved at $L = 24$). Only
the finite lower bounds are rigorous; the extrapolation is heuristic. These
deeper data leave the rigorous content untouched, but they show that the
bounded-$C(a)$ phenomenon of Conjecture~\ref{conj:correction} is robust to
the $L = 16$ truncation --- the band is wider but stays finite --- and they
give direct empirical support to $\alpha(a) < 2$
(Conjecture~\ref{conj:alpha}) across every tested multiplier.
\end{remark}

% =============================================================================
\section{Degenerate edge case: \texorpdfstring{$\lvert a \rvert = 1$}{|a| = 1}}\label{sec:edge}
% =============================================================================

\begin{theorem}[$a = +1$ vacuous]\label{thm:edge-case}
For $a = +1$ and every odd $b \ne 0$,
\[
\Obs_L^{(1,b)} = \emptyset \qquad \forall L \ge 1.
\]
\end{theorem}

\begin{proof}
The diophantine identity~\eqref{eq:diophantine} with $a = 1$ becomes
$1 = \sum_{j=0}^{J} X_j$, with $X_j = 2^{\Vk^{(j)}} - 2^{\Vi^{(j)} + v}$.
Since $\Vk^{(0)} = 0$, the $j = 0$ term contributes $2^{\Vk^{(0)}} = 1$;
moving it across, the identity rearranges to
\[
\sum_{j=1}^{J} 2^{\Vk^{(j)}} \;=\; \sum_{j=0}^{J} 2^{\Vi^{(j)} + v}.
\]
Each step adds a positive valuation ($\vk^{(i)} = \vtwo(a_K^{(i)} + b)
\ge 1$ because $a_K^{(i)}, b$ are odd), so the exponents
$\Vk^{(1)} < \Vk^{(2)} < \cdots < \Vk^{(J)}$ are strictly increasing, and
likewise $\Vi^{(0)} + v < \cdots < \Vi^{(J)} + v$. Hence the left-hand
side is a sum of $J$ \emph{distinct} powers of $2$ and the right-hand
side a sum of $J + 1$ distinct powers of $2$. By uniqueness of binary
representation, two equal positive integers have the same number of
binary ones; this forces $J = J + 1$, a contradiction. Therefore no
odd $r$ satisfies $\Phi_J = 0$ for $a = 1$, and $\Obs_L^{(1,b)} =
\emptyset$ for every $L$ and every odd $b$. (The conclusion is confirmed
by direct enumeration for all tested $b$ and $L \le 16$.) \qedhere
\end{proof}

\begin{remark}[The degenerate case $a = -1$]\label{rem:a-minus-1}
For $a = -1$ the situation is different: the map value $an + b = -n + b$
vanishes at $n = b$, where the parallel reduction degenerates
($\vtwo(0) = \infty$ forces immediate termination). The resulting
``obstructions'' are an artefact of this degeneration, not of genuine
synchronisation dynamics: direct enumeration gives
\[
\lvert \Obs_L^{(-1,b)} \rvert = 2^{L-2} \qquad (L \ge 2),
\]
exactly half the odd residues (for $b = +1$ they are the residues
$r \equiv 1 \pmod 4$). Thus $a = -1$ is \emph{not} vacuous; it is
degenerate. Neither value of $\lvert a \rvert = 1$ represents the
genuine obstruction theory, which requires $\lvert a \rvert \ge 3$: at
$a = +1$ the family is empty, and at $a = -1$ the map leaves
$\mathbb{N}_{\mathrm{odd}}$ through the zero of $an + b$. The order
lemma (Lemma~\ref{lem:order-lemma}) is therefore necessary but not
sufficient for a non-degenerate non-empty family.
\end{remark}

% =============================================================================
\section{Relation to the negative-Mersenne case}\label{sec:neg-mersenne}
% =============================================================================

For $a = -(2^{v_M} - 1)$ Mersenne with $v_M \ge 2$ (so $\lvert a \rvert
\in \{3, 7, 15, 31, 63, \ldots\}$), the parallel reduction admits a
finer structure than the general $(an+b)$ theory captures. The
algebraic identity
\[
-q r + b = 2^v (-q m + b) + (2^{v_M} - 2^v) b, \qquad q := \lvert a \rvert,
\]
has a vanishing correction term $(2^{v_M} - 2^v) b$ at $v = v_M$ that
synchronises the principal class of obstructions. Iterating produces a
sub-class hierarchy indexed by the signed shift index $j \in \Z$ (the
$A$-, $B$-, $C$-, $F_1$-, and $G_1$-families). The density bound
$\cWa \ge 1/(2 \lvert a \rvert)$ becomes asymptotically sharp as
$v_M \to \infty$; with $\lvert a \rvert = 2^{v_M} - 1$ and $e = v_M$ this
is $1/(2 \lvert a \rvert) = \tfrac12 \cdot 4^{-e/2}(1 + o(1)) =
\Theta(2^{-e})$. Note that this is the \emph{square root} of the
positive-$a$ scale: the negative-Mersenne density scales as
$\Theta(2^{-e})$, distinct from --- and much larger than --- the empirical
$\Theta(4^{-e})$ of the positive multipliers in Table~\ref{tab:scaling},
consistent with the cardinality ratio
$\lvert \Obs^{(-q)} \rvert / \lvert \Obs^{(+q)} \rvert \approx 2^{e}$
recorded in Appendix~\ref{sec:err-A3}. We treat this case in a separate
manuscript and only note the relationship.

% =============================================================================
\section*{Funding and competing interests}
% =============================================================================

The author received no funding for this research and declares no competing
interests.

\section*{Use of AI tools}

This manuscript was prepared with substantial assistance from generative AI
tools. Anthropic's Claude assisted with developing and formulating the proofs,
drafting and refining the exposition, and writing the accompanying verification
code; OpenAI's ChatGPT was used additionally for independent second-opinion
reviews of the arguments. All definitions, theorems, proofs, and results were
reviewed and independently checked by the author and are backed by the
reproducible verification and falsification code accompanying this manuscript. The
author takes full and sole responsibility for the correctness and content of
this manuscript. No AI system is listed as an author.

% =============================================================================
\appendix
\section{Refuted hypotheses}\label{sec:errata}
% =============================================================================

\subsection{Refuted hypothesis: inclusion \texorpdfstring{$\Obs^{(+q)} \subseteq \Obs^{(-q)}$}{O(+q) subset O(-q)}}\label{sec:err-A3}

An early conjecture (henceforth A3) asserted an injection
$\Obs^{(+q)} \hookrightarrow \Obs^{(-q)}$ with complement equal to a
Mersenne principal class. Concrete counterexamples (taken at $b = 1$, with
$\Obs^{(\pm q)} := \Obs_L^{(\pm q, 1)}$; unlike the $b$-independent
cardinalities, the specific witness residues depend on $b$ --- at $b = -1$
the same residues lie in $\Obs^{(-q)} \setminus \Obs^{(+q)}$ instead. The
first two appear already at level $L = 10$, the third at $L = 12$, since
for $\lvert a \rvert = 31$ the smallest non-vacuous level is $L = 11$):
\[
r = 29 \in \Obs^{(+3)} \setminus \Obs^{(-3)},
\quad
r = 9 \in \Obs^{(+7)} \setminus \Obs^{(-7)},
\quad
r = 33 \in \Obs^{(+31)} \setminus \Obs^{(-31)}.
\]
The empirical cardinality ratio is $\lvert \Obs^{(-q)} \rvert /
\lvert \Obs^{(+q)} \rvert \approx 2^e$ for small $e$ and approaches $2$
as $L \to \infty$, but with positive and negative obstructions
structurally distinct. The lesson: the sign of the multiplier is not
neutralisable via residue conjugation. A multiplicative conjugation
across the sign axis exists only as a partial bijection (revised
hypothesis A3$'$); concrete candidates such as $r \mapsto -r \bmod 2^L$
or $r \mapsto b r^{-1} \bmod 2^L$ remain empirically open.

\subsection{Refuted hypothesis: \texorpdfstring{$J_{\max}(L) = O(\log L)$}{Jmax(L) = O(log L)} for atoms}\label{sec:err-Jmax}

An earlier hypothesis based on the empirical mean $\bar J \sim \log_2 L
+ e$ asserted that $J_{\max}(L) = O(\log L)$ as well. This is refuted by
Theorem~\ref{thm:Lminus4-atom}: $J_{\max}(L, 3) = L - 4$ grows linearly.
The corrected picture is a bell-shaped distribution with logarithmic mode
and linear maximum --- a long-tail distribution. For the open problem
whether $\alpha(a) = 1$, it is the \emph{concentration} of the
distribution around its mode that matters, not a hard bound on $J_{\max}$.

% =============================================================================
\bibliographystyle{amsplain}
\bibliography{references}

\end{document}
